\subsection{A weaker criterion from the deck action on two-primary torsion}

\begin{claim}{A.C2.7}{Proved}
Let $C$ be a smooth proper geometrically connected curve of genus at
least two over a field $K$ finitely generated over $\F_2$. Suppose that,
over $\overline K$, it admits a connected nontrivial finite \'etale
Galois cover $D\to C$ of odd degree, with group $G$, for which
\[
             G\longrightarrow
             \operatorname{Aut}_{\F_2}(J_D[2](\overline K))
                                                               \tag{C2.19}
\]
is not faithful. Then the geometrically prime-to-two point-section
map of $C$ is injective. The conclusion also holds for its
finite-extension forms and hyperbolic open subcurves.

More precisely, a pair of distinct points giving equal section classes
would force (C2.19) to be faithful for every such cover. If
$m_2(G)$ denotes the least dimension of a faithful representation of
$G$ over $\F_2$, the same hypothetical pair would force
\[
                            f(D)-f(C)\ge m_2(G).               \tag{C2.20}
\]
In particular, a single such cover with $f(D)=f(C)$ suffices for
point-section injectivity.
\end{claim}
\begin{proof}
Work first over an algebraic closure. The curve $D$ is not
hyperelliptic. Otherwise its hyperelliptic involution is unique,
so every deck transformation descends to its genus-zero quotient.
Choose an element $\sigma\in G$ of odd prime order. Its image on
$\mathbf P^1$ has the same order, since the kernel of descent is
the order-two hyperelliptic group. A finite odd-order automorphism
of $\mathbf P^1$ in characteristic two has a fixed geometric
point. Its fibre on $D$ has at most two geometric points, so
$\sigma$ fixes a point in that fibre. This contradicts the free
deck action of an \'etale Galois cover.

Let $P',Q'\in D(\overline K)$ lie over distinct points of $C$.
The class $T=[Q'-P']$ has trivial stabilizer in $G$. Indeed,
$\delta T=T$ for $1\ne\delta\in G$ would give
\[
                     \delta Q'+P'\sim Q'+\delta P'.
                                                               \tag{C2.21}
\]
The two effective degree-two divisors in (C2.21) are distinct:
their terms lying over each of the two base points would
otherwise force both $\delta P'=P'$ and $\delta Q'=Q'$.
A linear equivalence between distinct effective degree-two
divisors on a curve of genus at least two gives a degree-two
map to $\mathbf P^1$, after removing any common base point.
A degree-one map is excluded by the genus. Thus (C2.21) would
make $D$ hyperelliptic, a contradiction.

The \'etale two-adic Tate module of $J_D$ is free of rank $f(D)$,
and its geometric two-primary torsion is the quotient of its
rational Tate module by this lattice. An odd-order automorphism
acting trivially on $J_D[2](\overline K)$ acts trivially on the
whole two-adic module. This can be checked without any
torsion-free congruence-group citation: if a matrix $A$ of odd
order $m$ satisfies $A\equiv1\pmod{2^r}$, then
\[
 A^m\equiv1+m(A-1)\pmod{2^{r+1}}
\]
implies $A\equiv1\pmod{2^{r+1}}$. Starting at $r=1$ and
iterating proves $A=1$. Thus any element in the kernel of
(C2.19) fixes every two-primary torsion point of $J_D$.

Now suppose that distinct $P,Q\in C(K)$ give equal section
classes. Descend the cover and its deck transformations after
a finite extension of $K$, and extend again to make a chosen
lift $P'$ rational. By A.3 there is a lift $Q'$ with
$T=[Q'-P']$ two-primary. A nonidentity element in the kernel
of (C2.19) would fix $T$, contrary to its trivial stabilizer.
This proves the criterion and the necessity of faithfulness.
The same argument after restriction proves the stated form and
open-subcurve variants.

For the quantitative assertion, the averaging idempotent
$|G|^{-1}\sum_{g\in G}g$ acts integrally on the two-adic Tate
module, since $|G|$ is odd. Its invariant summand is the
pullback of the two-adic Tate module of $J_C$: pullback
followed by norm is multiplication by $|G|$, and norm followed
by pullback is the sum of deck transformations. This summand
has rank $f(C)$. The complementary summand has rank
$f(D)-f(C)$, and its reduction modulo two must be a faithful
$G$-representation because the invariant summand is trivial.
This proves (C2.20).
\end{proof}

\begin{claim}{A.C2.8}{Proved}
In the genus-two cyclic situation of A.C2.5, let $S$ be the
union of the nonzero doubling orbits that miss $\Theta$.
The kernel of the deck action on geometric two-primary torsion
has order
\[
                    d=\gcd\bigl(n,\{i:i\in S\}\bigr),        \tag{C2.22}
\]
where $d=n$ if $S$ is empty. Consequently $d>1$ already implies
point-section injectivity by A.C2.7.

In particular, if $n=\ell^a$ is an odd prime power, it suffices
that every doubling orbit on $(\mathbf Z/n)^\times$ meets
$\{i:L^i\in\Theta\}$. No condition is needed on characters
of smaller order. The simple sufficient condition is
\[
 L\in\Theta\text{ of exact order }\ell^a,\qquad
                 \langle2,-1\rangle=(\mathbf Z/\ell^a)^\times.
                                                               \tag{C2.23}
\]
For a cyclic group of odd prime-power order $\ell^a$ one has
$m_2(G)=\operatorname{ord}_{\ell^a}(2)$.
\end{claim}
\begin{proof}
The Artin--Schreier description of \'etale $H^1$ identifies its
scalar extension with the stable Frobenius part of
$H^1(D,\mathcal O_D)$. The calculation of A.C2.5 therefore
shows that its nontrivial cyclic characters are exactly those
indexed by $S$, with multiplicity one. Dualizing to geometric
two-torsion may invert all characters; it does not change
their common kernel. A power $\delta^k$ acts trivially exactly
when $n\mid ik$ for every $i\in S$. The number of such powers
is the greatest common divisor in (C2.22).

If $n=\ell^a$, this kernel is nontrivial exactly when no
primitive character survives, since a primitive index is
coprime to $n$, whereas all nonprimitive indices are divisible
by $\ell$. Symmetry of $\Theta$ gives (C2.23) as in A.C2.5.
For example at $n=9$, a theta point of exact order nine kills
the single orbit of all six primitive characters. The
order-three block may remain ordinary, but the deck action
then has a kernel of order three, which suffices here. This
does not assert existence of such a nine-torsion theta point
on every curve.

Finally, an irreducible $\F_2$-representation of a cyclic
group corresponds to a doubling orbit of its geometric
characters. A faithful character of order $\ell^a$ has orbit
length $\operatorname{ord}_{\ell^a}(2)$. Any faithful
representation of this cyclic prime-power group must contain
such a character: otherwise all character orders divide
$\ell^{a-1}$ and a subgroup of order $\ell$ acts trivially.
One faithful orbit already gives a faithful representation,
proving the formula for $m_2(G)$.
\end{proof}

The full-quotient obstruction in A.C2.6 and the weaker deck
criterion are compatible. Surviving order-three characters
prevent the entire Prym from having $2$-rank zero, but at
degree $3^a$, $a>1$, they need not make the deck action
faithful. Torsion density still does not guarantee the
prime-power order or simultaneous primitive-orbit
intersections required by (C2.23). A degree-sized lower
bound in place of (C2.20), or a construction forcing a
sufficiently small normalized $2$-rank in general towers,
has not been established here.

\subsection{What the entire first cubic tower still cannot detect}

\begin{claim}{A.C2.9}{Proved; an obstruction to the specified target method}
Let $C/\F_4$ be the ordinary genus-two curve
\[
 y^2+y=x+1/x+\alpha/(x+1),\qquad \alpha^2+\alpha+1=0.
\]
Its full mod-three abelian \'etale cover $U\to C$ has group
$(\Z/3)^4$, degree $81$, genus $82$, and ordinary Jacobian.
Every intermediate curve is ordinary. Thus neither $U$ nor an
intermediate curve admits a nonconstant map to an abelian variety
of two-rank zero. The deck action on geometric two-primary torsion
is faithful for every nontrivial intermediate Galois cover of $C$.
The quotient of $U$ by a lifted hyperelliptic involution has genus
$40$ and is ordinary. These statements concern the complete first
cubic compositum; they do not imply the corresponding assertions
for higher three-power covers and do not refute A1.
\end{claim}
\begin{proof}
Work geometrically, and base the Abel map at a Weierstrass point.
Its pullback of $[3]:J_C\to J_C$ is connected because the Abel
map induces the maximal abelian quotient of the geometric
fundamental group \citep[III, Proposition 9.1]{A_Milne2008}.
Hence $G=\Gal(U/C)=(\Z/3)^4$ and Riemann--Hurwitz gives
$g(U)=1+81=82$. There are $40$ index-three subgroups $K$ of
$G$, and $D_K=U/K$ runs through all the connected cyclic cubic
covers. Write $P_K$ for the connected Prym of $D_K\to C$.
Claim A.C2.2 proves
\[
  J_{D_K}\sim J_C\times P_K,\qquad P_K\sim E_K^2,
\]
where $E_K$ is an ordinary elliptic reflection quotient and
$\dim P_K=2$.

Here is an exact exhaustion of the Jacobian factors. For a subgroup
$H\subset G$ put $e_H=|H|^{-1}\sum_{h\in H}h$ in $\Q[G]$.
The idempotents $e_G$ and $e_K-e_G$, for the $40$ hyperplanes
$K$, are pairwise orthogonal. For distinct $K,L$, one has
$e_Ke_L=e_G$ because $KL=G$. Every nonzero vector of $\F_3^4$
lies in $13$ hyperplanes, whereas zero lies in $40$, so
\[
 \sum_K\sum_{h\in K}h=27\,1+13\sum_{g\in G}g,
 \qquad e_G+\sum_K(e_K-e_G)=1.
\]
Pullback and norm identify these factors with $J_C$ and the
pullbacks of the $P_K$. Their sum therefore gives an isogeny
\[
                  J_U\sim J_C\times\prod_{40}E_K^2.       \tag{C2.24}
\]
All denominators in the displayed idempotents are powers of three,
so they split the two-adic Tate module integrally as well; no
two-primary component is omitted. The factor dimensions total
$2+40\cdot2=82$.

Ordinariness is preserved under products and isogenies, so $U$ is
ordinary. Pullback and norm make the Jacobian of every intermediate
curve an isogeny factor of $J_U$, hence ordinary. The image of a
homomorphism from an ordinary abelian variety to one of two-rank
zero has dimension zero: it is both an ordinary quotient and a
two-rank-zero subvariety. The Albanese property proves the assertion
about maps of curves.

Each nontrivial character pair of $G$ occurs in exactly one $P_K$.
Its kernel is $K$. Since $\bigcap_K K=1$, these characters act
faithfully on geometric two-torsion. For an intermediate cover
$U/H\to C$, every nontrivial character of $G/H$ still occurs,
with the same conclusion. Thus the sufficient criterion A.C2.7
does not apply to these covers.

The involution lifts on the Abel pullback as
$(P,z)\mapsto(\iota P,-z)$, and inverts $G$. On the fibre over
each of the three hyperelliptic branch points it is an affine
inversion of the odd $G$-torsor, with exactly one fixed point.
The local \'etale map to $C$ preserves its different exponent two.
The total different is six, whence
\[
 162=2\bigl(2g(U/\langle\tau\rangle)-2\bigr)+6
\]
and the quotient genus is $40$. Its Jacobian is again an ordinary
isogeny factor of $J_U$.

The finite integer certificate
\path{code/A_char2_cubic_tower_certificate.py} verifies all $40$
hyperplanes, their incidence counts, $861$ group-ring convolution
identities, the character partition, and the unique fixed point of
every affine inversion. The proof above supplies the identification
of these group-ring calculations with actual Jacobian factors.
\end{proof}

\begin{claim}{A.C2.10}{Proved; the inertia obstruction in equivariant odd towers}
Let $C$ be any genus-two curve over an algebraically closed field
of characteristic two. Let $V\to C$ be an odd-degree \'etale
Galois cover equivariant under the hyperelliptic involution, so
that $V\to\mathbf P^1$ has Galois group $\Gamma$ in an exact
sequence
\[
             1\longrightarrow N\longrightarrow\Gamma
               \longrightarrow C_2\longrightarrow1.
\]
Every map between intermediate curves of this $\Gamma$-cover
has ramification indices one or two. No odd-order subgroup of
$\Gamma$ has a genus-one quotient of $V$. Consequently the
usual elliptic group quotients in this tower cannot supply the
odd common ramification index in A.21 or a prime-to-two Galois
elliptic quotient for A.29l. This statement does not exclude maps
or additional automorphisms outside the specified tower group.
\end{claim}
\begin{proof}
Since $V\to C$ is \'etale, the ramification indices of
$V\to\mathbf P^1$ are those of $C\to\mathbf P^1$, namely
one and two. If $I_R$ is inertia at a point of $V$, the inertia
of $V\to V/H$ is $I_R\cap H$ for any subgroup $H\subset\Gamma$.
All intermediate ramification indices therefore divide two.

An odd-order subgroup $H$ maps trivially to $C_2$, so lies in
$N$. The cover $V/H\to C$ is \'etale, of degree $[N:H]$, and
\[
                 g(V/H)=1+[N:H](g(C)-1)\ge2.
\]
This proves both obstructions. Equivariant refinements are available
for any given odd Galois cover: intersect its normal subgroup of
the geometric fundamental group with the transform under the
hyperelliptic involution. The resulting odd quotient is invariant.
One can realize a lifted involution by normalizing the fibre product
with the conjugate cover and taking a component fixed by factor
interchange; the number of components is odd. Thus the obstruction
persists after these refinements, without asserting anything about
arbitrary maps outside their group quotients.
\end{proof}

\subsection{An exact nine-torsion point resolves the cubic-tower example}

\begin{claim}{A.C2.11}{Proved}
The curve
\[
 C/\F_4:\quad y^2+y=x+1/x+\alpha/(x+1),
 \qquad\alpha^2+\alpha+1=0,
\]
has an injective geometrically prime-to-two point-section map over
every finite extension. The conclusion also holds for its
finite-extension forms over fields finitely generated over $\F_2$
and for its hyperbolic open subcurves. An explicit cyclic \'etale
degree-nine cover has genus ten, two-rank four, and an order-three
deck subgroup acting trivially on all geometric two-primary torsion.
This higher cover overcomes the valid first-cubic obstruction
A.C2.9; the entire Prym is not required to have two-rank zero.
\end{claim}
\begin{proof}
Let $W_0,W_1,W_\infty$ be the unique points above $0,1,\infty$.
The different exponent at each simple Artin--Schreier pole is two,
so
\[
 \operatorname{div}(dx)=2W_0+2W_1-2W_\infty,
 \qquad B=W_0+W_1-W_\infty.
\]
The line bundle of locally exact differentials is $\mathcal O(B)$:
locally its transition functions square to those of ordinary
differentials, and the exact differential $dx$ has the displayed
divisor. The Raynaud theta is consequently
\[
                 \Theta=\{\mathcal O(R-B):R\in C(\overline\F_2)\}.
\]

Choose $t$ with $t^4+t^3+t^2+t+1=0$ and put
$\alpha=t+t^{-1}=t^3+t^2+1$. Then $\alpha^2+\alpha+1=0$.
Choose $\beta$ with $\beta^2+\beta=t+1$.
The identity
\[
                 t+t^{-1}+\alpha/(t+1)=t+1
\]
makes $R=(t,\beta)$ a point of $C$. These choices lie in
$\F_{16}$ and $\F_{256}$ respectively; the absolute trace of
$t+1$ from $\F_{16}$ is one.

Define
\[
\begin{aligned}
 b(x)&=x^3+(t^2+t+1)x^2+(t^2+t+1)x+t^2,\\
 a(x)&=x^3+t x^2+t x+t^2,\\
 h&=a(x)+b(x)(y+\beta),\qquad u=x(x+1).
\end{aligned}
\]
Exact polynomial multiplication modulo $t^4+t^3+t^2+t+1$ gives
\[
 \begin{aligned}
 u(a^2+ab)+b^2\bigl(x^3+x^2+(1+\alpha)x+1+(t+1)u\bigr)
                          &=(x+t)^9,\\
 \operatorname{Nm}_{C/\mathbf P^1}(h)&=(x+t)^9/u.
 \end{aligned}                                                   \tag{C2.25}
\]
Here $a(t)=0$ and $b(t)=1$, so $h(R)=0$ and its conjugate
value at $R$ is one. Thus all nine zeros lie at $R$.
The poles are exactly simple at $W_0,W_1$ and of order seven
at $W_\infty$: $b(0)=t^2$ and $b(1)=1+t^2$ are nonzero,
and $x^3y$ has the unique largest pole order seven at infinity.
Therefore
\[
 \operatorname{div}(h)=9R-W_0-W_1-7W_\infty,
 \qquad \operatorname{div}(h/u^4)=9(R-B).                         \tag{C2.26}
\]
The second equality uses
$\operatorname{div}(u)=2W_0+2W_1-4W_\infty$.

Hence $\eta=\mathcal O(R-B)$ lies on $\Theta$ and is killed
by nine. It is not killed by three. Claim A.C2.2 and the exact
theta formula A.C2.5 give $\Theta\cap J_C[3]=\varnothing$:
every cubic Prym is ordinary since
$(1+1+\alpha)^3+\alpha=1+\alpha\ne0$, and zero is absent
from $\Theta$ because $C$ is ordinary. Thus $\eta$ has exact
order nine.

Let $D\to C$ be the connected cyclic \'etale degree-nine cover
associated with $\eta$, obtained by taking a ninth root of $h/u^4$.
Its genus is ten. The primitive characters form the single doubling
orbit $1,2,4,8,7,5$ modulo nine. Since $\eta\in\Theta$, Frobenius
is nilpotent on their six one-dimensional cohomology spaces.
The orbit $3,6$ survives, since there is no nonzero three-torsion
on $\Theta$. Formula (C2.16) gives $f(D)=2+2=4$.
The common kernel of the surviving nontrivial characters $3,6$
is the order-three subgroup of the deck group. By A.C2.8 it
acts trivially on every geometric two-primary torsion point.

A section collision on $C$ lifts by A.3 to a two-primary difference
$[Q'-P']$ on $D$. A nonidentity element $\sigma$ of this
order-three kernel fixes that difference, giving
\[
                     Q'+\sigma P'\sim P'+\sigma Q'.
\]
The effective degree-two divisors are distinct, since the lifts
lie over distinct base points and the deck action is free. This
would make $D$ hyperelliptic, contrary to A.C2.7. This proves A1 for
the stated curve.
All data descend after a finite extension; the same lifting and
Kummer argument proves the finitely generated and open-subcurve
variants.

The exact certificate \path{code/A_theta_nine_certificate.py}
checks (C2.25), every local value, the trace and the two character
orbits. Root independently verified the polynomial identity with a
separate quotient-polynomial implementation in
\path{code/A_theta_nine_independent.py}; its full mathematical
review is \path{audit/A_theta_nine_root_review.md}.
\end{proof}

\subsection{Prime-power theta torsion is confined to a finite division field}

\begin{claim}{A.C2.12}{Proved; a finite search bound, not an existence theorem}
Let $J/\F_q$ be a genus-two Jacobian and let $\Theta$ be a
geometrically irreducible theta divisor defined over $\F_q$.
Let $\ell$ be an odd prime different from the characteristic,
and choose $r$ such that $J[\ell]$ is rational over $\F_{q^r}$.
Then
\[
 \Theta(\overline\F_q)\cap J[\ell^\infty]
       \subset J(\F_{q^r})[\ell^\infty].                         \tag{C2.27}
\]
In particular the intersection is finite. For the Raynaud theta
of the curve in A.C2.11, every three-primary point is killed by
$27$. Since theta three-torsion is absent, its only possible
nonzero three-primary orders are nine and $27$; existence of
order nine was proved separately, and order $27$ is not inferred.
\end{claim}
\begin{proof}
This is a direct genus-two form of the bounded-prime-support
phenomenon for almost rational torsion
\citep[Proposition 2, p.16]{A_BoxallGrant2006}. We give the
fixed-field proof, including the case $\ell=3$.

Write $\pi$ for $q$-power Frobenius and $\sigma=\pi^r$.
On $T_\ell J$, $\sigma$ is congruent to one modulo $\ell$.
Suppose $\eta\in\Theta\cap J[\ell^\infty]$ is not fixed.
The principal congruence subgroup modulo $\ell^a$ is an
$\ell$-group, so there is a least $b\ge1$ with
$\sigma^{\ell^b}\eta=\eta$. Put
$\tau=\sigma^{\ell^{b-1}}=1+\ell M$ on $T_\ell J$.
The operator
\[
 U=\frac{1+\tau+\cdots+\tau^{\ell-1}}\ell
   =1+\sum_{j=1}^{\ell-1}
       \binom{\ell}{j+1}\ell^{j-1}M^j
\]
is integral and congruent to one modulo $\ell$. At $\ell=3$
the final coefficient is three, so this case is included.
Thus $U$ acts invertibly on geometric $\ell$-primary torsion
and commutes with $\tau$. If $T=\tau\eta-\eta$, then
$T\ne0$ and
\[
                 0=(\tau^\ell-1)\eta=\ell U T,
                 \qquad \ell T=0.
\]
Since $\tau$ fixes $J[\ell]$, it fixes $T$. Its orbit on
$\eta$ is therefore the full coset
$\eta+\langle T\rangle$, of size $\ell$. Every point of
that coset lies on $\Theta$, because $\tau$ is a Frobenius
power, and also on $\Theta-T$.

The two irreducible divisors $\Theta$ and $\Theta-T$ are
distinct: a translation preserving a principal theta divisor
lies in the trivial kernel of its principal polarization.
Their intersection number is $\Theta^2=2$, so they cannot
contain $\ell\ge3$ distinct common points. This contradiction
proves $\sigma\eta=\eta$, which is (C2.27). No scalar or
semisimple Frobenius hypothesis was used.

For the curve in A.C2.11, direct counts give
$\#C(\F_4)=5$ and $\#C(\F_{16})=15$. The Frobenius
polynomial is $P(X)=X^4-X^2+16$. Exact reduction gives
\[
          X^{12}-1\equiv3(235X^2-1115)\pmod{P(X)},
\]
so $\pi^{12}$ fixes $J[3]$. By (C2.27) every three-primary
theta point is fixed by $\pi^{12}$. The further exact identity
\[
\begin{aligned}
 &(X^{12}-1)(176+47X^2)+27\cdot41440\\
 &\quad=P(X)\bigl(47X^{10}+223X^8-529X^6-4097X^4
                              +4367X^2+69919\bigr)
\end{aligned}                                                     \tag{C2.28}
\]
then kills it by $27\cdot41440$. Since $41440$ is prime to
three, its order divides $27$. These are polynomial
endomorphism identities; nontrivial Jordan reduction causes
no additional assumption. The counts and both identities are
verified in \path{code/A_theta_frobenius_certificate.py}.
Root independently checked the proof and integer remainders
in \path{audit/A_theta_prime_power_root_review.md}.
\end{proof}

Thus repeated division does not force theta points of arbitrarily
high prescribed prime-power order. For each fixed prime, A.C2.12
reduces that route to a finite set. Neither prime-to-two torsion
density nor surjectivity onto selected primary components forces
all other primary components to vanish. The required existence
statement for arbitrary genus-two curves remains open here.

\subsection{Varying the odd prime: a full three-primary obstruction and an explicit seven-cover}

\begin{claim}{A.C2.13}{Proved}
Let $\alpha^2+\alpha+1=0$ and let
\[
 C/\F_4:\qquad y^2+y=x+\frac{\alpha}{x}
                         +\frac{\alpha^2}{x+1}.
\]
Its Raynaud theta contains no point of three-power order, including
the origin. Every geometrically connected finite \'etale abelian
three-group cover of $C$ is ordinary, and every nontrivial such
geometric deck group acts faithfully on two-primary torsion.
None of those curves admits a nonconstant map to a positive-dimensional
abelian variety of two-rank zero.

Nevertheless, theta contains the exact order-seven point below.
Its cyclic \'etale cover has genus eight and two-rank two, with
trivial deck action on geometric two-primary torsion. Consequently
A1 holds for this $C$ over every finite extension, for every
finitely generated form of this geometric curve, and for its
hyperbolic open subcurves.
\end{claim}
\begin{proof}
Let $W_0,W_1,W_\infty$ denote the points over $0,1,\infty$.
The three Artin--Schreier poles are simple, with different exponent
two. Thus $g(C)=2$ and
\[
 \operatorname{div}(dx)=2W_0+2W_1-2W_\infty.
\]
The regular independent logarithmic forms $dx/x$ and $dx/(x+1)$
have respective valuations $(0,2,0)$ and $(2,0,0)$ at these
three points. They are Cartier-fixed, so $C$ is ordinary.
As in A.C2.11, the locally exact differential line bundle is
\[
 \mathcal B_C=\mathcal O_C(B),\qquad B=W_0+W_1-W_\infty.
\]
For the convention
$L\in\Theta\Longleftrightarrow H^0(C,\mathcal B_C\otimes L)\ne0$,
its Raynaud theta is
\[
 \Theta=\{\mathcal O_C(R-B):R\in C(\overline\F_2)\}.
\]
It is geometrically irreducible, defined over $\F_4$, and
symmetric because $2B$ is canonical.

The cubic obstruction in A.C2.2 equals
\[
                   (1+\alpha+\alpha^2)^3+\alpha\alpha^2=1.
\]
Every geometrically connected cyclic cubic cover is therefore
ordinary. By the all-character criterion A.C2.5,
\[
                         \Theta\cap J_C[3]=\varnothing.          \tag{C2.29}
\]
The origin is absent too, since $C$ is ordinary.

\proofparagraph{The complete three-primary packet.}
Exact counts, including the three ramification points separately, are
\[
                \#C(\F_4)=3,\qquad \#C(\F_{16})=15.
\]
The curve trace formula and Newton identities give
\[
                       P(X)=X^4-2X^3+X^2-8X+16                 \tag{C2.30}
\]
for the Jacobian Frobenius polynomial
\citep[III, Theorem 11.1 and Corollary 11.4; II, Theorem 1.1]
{A_Milne2008}. The exact integer identities are
\[
             X^5-1\equiv3X^3+6X^2-33\pmod{P(X)}
\]
and
\[
 \begin{aligned}
 &(X^5-1)(-207X^3-57X^2+248X+856)-40344\\
 &\qquad=P(X)(-207X^4-471X^3-487X^2-1303X-2575).
 \end{aligned}                                                   \tag{C2.31}
\]
Write $\pi$ for $\F_4$-Frobenius. The first identity makes
$\pi^5$ trivial on $J[3]$. By A.C2.12, every three-primary
$\eta\in\Theta$ is fixed by $\pi^5$. The second identity and
$P(\pi)=0$ give $40344\eta=0$. Since
$40344=3\cdot13448$ with $3\nmid13448$, this implies $3\eta=0$.
Equation (C2.29) proves
\[
                      \Theta\cap J[3^\infty]=\varnothing.        \tag{C2.32}
\]
These endomorphism identities apply to the entire torsion group;
no scalar or semisimple reduction of Frobenius is assumed.

Let $D\to C$ be a geometrically connected finite \'etale abelian
three-group cover. Over the algebraic closure, split its characters
and decompose its direct-image algebra into torsion lines $L_\chi$.
Geometric connectedness makes every nontrivial character line
nontrivial. Each has three-power order and hence avoids theta.
Each $H^1(C,L_\chi)$ has dimension one. The locally exact
differential sequence therefore gives a bijective Frobenius arrow
from every nontrivial character space to its doubled-character
space. The invariant $H^1(C,\mathcal O_C)$ is ordinary as well.
Thus all of $H^1(D,\mathcal O_D)$ has bijective Frobenius, and
$D$ is ordinary.

Every nontrivial character occurs in the Frobenius-stable
cohomology and, dually, in the geometric two-primary
representation; duality merely inverts the characters.
Characters separate a finite abelian odd-order group, proving
faithfulness. Quotients of ordinary abelian varieties are
ordinary. A nonconstant map from $D$ to an abelian variety
generates a positive-dimensional quotient of $J_D$, so it
cannot land nontrivially in an abelian variety of two-rank zero.

\proofparagraph{The exact order-seven divisor.}
Choose $\beta^2+\beta=\alpha$ and put $R=(\alpha^2,\beta)$.
The curve equation at $x=\alpha^2$ has right side $\alpha$.
The point lies over $\F_{16}$; one may take
$\beta^4+\beta+1=0$ and $\alpha=\beta^2+\beta$.
Define
\[
 \begin{aligned}
 u&=x(x+1),& a&=x+\alpha^2,& b&=x^2+\alpha^2,\\
 h&=a+b(y+\beta),&
 F&=x^3+x^2+x+\alpha
     =u\left(x+\alpha/x+\alpha^2/(x+1)\right).
 \end{aligned}
\]
Exact polynomial multiplication over $\F_4$ gives
\[
          u(a^2+ab)+b^2(F+\alpha u)=(x+\alpha^2)^7,
 \qquad
          \operatorname{Nm}_{C/\mathbf P^1}(h)
                    =\frac{(x+\alpha^2)^7}{u}.                  \tag{C2.33}
\]
At $R$, $a(\alpha^2)=0$ and $b(\alpha^2)=1$, so $h$ vanishes
and its Artin--Schreier conjugate equals one. Since the $x$-map
is \'etale there, its sole zero is $R$, of exact order seven.
At $W_0,W_1$, the values of $b$ are $\alpha^2,\alpha$,
giving a simple pole at each. At infinity $x^2y$ is the unique
term of largest pole order five. Therefore
\[
 \operatorname{div}(h)=7R-W_0-W_1-5W_\infty,\qquad
 \operatorname{div}(h/u^3)=7(R-B),                               \tag{C2.34}
\]
using $\operatorname{div}(u)=2W_0+2W_1-4W_\infty$.

Thus $\eta=\mathcal O(R-B)$ lies on theta and is killed by seven.
It is nonzero: $R\sim B$ would make $2R$ canonical, forcing
$R$ to be a ramification point of the genus-two hyperelliptic
map, whereas $x(R)=\alpha^2\notin\{0,1,\infty\}$.
Its order is exactly seven. The normalization of
\[
                              z^7=h/u^3
\]
is consequently a geometrically connected \'etale cover
after a finite constant extension. Exact order proves
connectedness; divisibility of every valuation by seven proves
\'etaleness. Its genus is eight.

\proofparagraph{The successful deck kernel.}
The nonzero characters have the two doubling orbits
\[
                    \{1,2,4\},\qquad
                    \{3,6,5\}=-\{1,2,4\}.
\]
The first meets theta at $\eta$ and the second at $-\eta$.
Every character space has dimension one, so each cycle
contains a zero Frobenius arrow and is nilpotent.
Only the invariant rank-two ordinary part survives.
The cover has two-rank two and its full order-seven deck
group acts trivially on geometric two-primary torsion.

Claim A.C2.7 proves A1 for the stated curve. A collision would
lift to a two-primary point-divisor difference fixed by a
nonidentity deck transformation, producing a linear
equivalence between distinct effective divisors of degree two.
The upstairs curve is nonhyperelliptic: a free odd-order
automorphism on a hyperelliptic curve would descend to the
line, have a fixed point there, and hence fix a point in
a fibre of at most two points. The degree-two equivalence
is impossible. The argument persists under every finite
extension. Finite-base-field extension, A.8 spreading and
the proper-to-open reduction give the remaining stated scope.

The self-contained verifier is
\path{code/A_theta_seven_certificate.py}, with adjacent JSON.
Independent implementations are
\path{code/A_theta_seven_independent.py}, using a quadratic
field tower and companion matrices, and
\path{code/A_theta_seven_root_independent.py}, using separate
quotient-polynomial and enumeration calculations.
The complete independent geometric review is
\path{audit/A_theta_seven_independent_review.md}.
All three certificates pass; they verify the displayed
identities, field counts, local values and character orbits.
The divisor and cohomology arguments above supply their
mathematical implications.

This requires a change of odd prime for the specified
abelian theta/deck-kernel construction. It does not exclude
nonabelian three-group covers or decide A1 for arbitrary
ordinary genus-two curves.
\end{proof}

\subsection{No fixed finite set of odd primes suffices for the abelian theta method}

\begin{claim}{A.C2.14}{Proved; a method obstruction, not a counterexample to A1}
Let $J/\F_q$ be a genus-two Jacobian with a geometrically irreducible
theta divisor $\Theta/\F_q$. Let $S$ be a finite nonempty set of odd
primes different from the characteristic, and put $L=\prod_{\ell\in S}\ell$.
If a Frobenius power $\sigma$ fixes $J[L]$, it fixes every point
of $\Theta$ whose order is supported on $S$.

For every finite nonempty set $S$ of odd primes there exists an
ordinary smooth proper genus-two curve $C$ over a finite field of
characteristic two such that
\[
                    \Theta_C\cap J_C[S^\infty]=\varnothing.
                                                                    \tag{C2.35}
\]
Every geometrically connected finite \'etale cover that is abelian
over this original $C$, with group order supported on $S$, is
ordinary and has faithful geometric deck action on two-primary
torsion. None admits a nonconstant map to a positive-dimensional
abelian variety of two-rank zero. The curve and its finite field
may depend on $S$. No assertion is made about arbitrary iterated
abelian covers whose composite group is nonabelian.
\end{claim}
\begin{proof}
\proofparagraph{A mixed-primary fixed-field lemma.}
Let $\eta\in\Theta$ have order supported on $S$. On
$J[\ell^a]$, the operator $\sigma$ belongs to the principal
congruence kernel modulo $\ell$, an $\ell$-group.
Thus the length $n$ of its orbit on $\eta$ is supported on $S$.
If $n>1$, choose $\ell\mid n$ and put $\tau=\sigma^{n/\ell}$.
Its orbit on $\eta$ has length $\ell$. For $r\in S$ distinct
from $\ell$, its orbit on the $r$-primary component has both
$r$-power order and order dividing $\ell$, so that component
is fixed. Therefore $T=\tau\eta-\eta$ is nonzero and
$\ell$-primary.

On $T_\ell J$, write $\tau=1+\ell M$. The integral operator
\[
 U=\frac{1+\tau+\cdots+\tau^{\ell-1}}\ell
   =1+\sum_{j=1}^{\ell-1}
       \binom{\ell}{j+1}\ell^{j-1}M^j
\]
is congruent to one modulo $\ell$, including at $\ell=3$.
Thus $\ell UT=(\tau^\ell-1)\eta=0$ implies $\ell T=0$.
Since $\tau$ fixes $J[\ell]$, it fixes $T$. The entire coset
$\eta+\langle T\rangle$, of size $\ell$, lies in both
$\Theta$ and $\Theta-T$. These irreducible divisors are
distinct because a principal polarization has trivial
translation stabilizer, and their intersection number is two.
This contradicts $\ell\ge3$ and proves $\sigma\eta=\eta$.

In particular, if actual Frobenius $\pi$ satisfies
\[
 \pi-1=\ell A_\ell,\qquad
 A_\ell\in\operatorname{GL}(T_\ell J)\quad(\ell\in S),             \tag{C2.36}
\]
then its kernel on $J[\ell^\infty]$ is exactly $J[\ell]$.
The fixed-field lemma yields
\[
            \Theta\cap J[S^\infty]\subseteq\Theta\cap J[L].
                                                                    \tag{C2.37}
\]
This treats arbitrary primary mixtures and does not assume
scalar or semisimple Frobenius.

\proofparagraph{A nonempty smooth locus avoiding the fixed level.}
In ordinary genus-two moduli in characteristic two, the
condition $\Theta_C\cap J_C[L]=\varnothing$ is open.
On the relative Frobenius-twisted smooth curve, the locally
exact differential sheaf is a line bundle. The condition
$H^0(\mathcal B\otimes M)\ne0$ is closed by upper
semicontinuity on the finite \'etale $L$-torsion parameter
space; its finite image in moduli is closed. The origin is
avoided on the ordinary locus.

Here is a complete finite-level degeneration proof that this
open is nonempty. Over an algebraic closure $\kappa$ of $\F_2$,
join two ordinary elliptic curves to form $X_0=E_1\cup E_2$
with one node. The stable deformation calculation of
\citet[Proposition 1.5, its algebraization paragraph, and
Theorem 1.6, pp.81--83]{A_DeligneMumford1969} gives a projective
regular smoothing $f:X\to\operatorname{Spec}R$, where
$R=\kappa[[t]]$, with local node equation $uv=t$ and smooth
genus-two generic fibre. A smooth special point lifts to a
section, if needed to rigidify line bundles.

The normalization sequence for units gives
$\operatorname{Pic}^0(X_0)=\operatorname{Pic}^0(E_1)\times
\operatorname{Pic}^0(E_2)$, since its dual graph is a tree.
There are exactly $L^4$ line-bundle classes killed by $L$.
A trivialization of the $L$th power of each gives a finite
\'etale $\mu_L$-torsor. Proper henselian equivalence extends
each torsor to $X$
\citep[Lemma 58.9.1, tag 0A48]{A_StacksProperHenselian}.
Full faithfulness extends its action and the torsor identity.
Its character line $\mathcal M$ therefore satisfies
$\mathcal M^L\simeq\mathcal O_X$ as an actual line bundle.

These $L^4$ classes remain distinct on the geometric generic
fibre. Indeed, a quotient line becoming geometrically trivial
is already trivial over the generic field by proper cohomology
and field base change. A generic trivialization on the
regular surface has vertical divisor, so the line is
$\mathcal O_X(aE_1+bE_2)$. Since $E_1E_2=1$ and
$E_i^2=-1$, its two multidegrees are $(b-a,a-b)$.
Torsion forces them to vanish. Thus $a=b$, and the line
is $\mathcal O_X(aX_0)=\mathcal O_X$ because the total
fibre is principal. Its special class was trivial as well.
The $L^4$ distinct generic classes exhaust the smooth
genus-two Jacobian's $L$-torsion.

Write $\mathcal M_i=\mathcal M|_{E_i}$. Normalization gives
\[
 \begin{aligned}
 0&\longrightarrow H^0(X_0,\mathcal M)
 \longrightarrow H^0(E_1,\mathcal M_1)\oplus H^0(E_2,\mathcal M_2)
 \longrightarrow\mathcal M_{\rm node}\\
 &\longrightarrow H^1(X_0,\mathcal M)
 \longrightarrow H^1(E_1,\mathcal M_1)\oplus H^1(E_2,\mathcal M_2)
 \longrightarrow0 .
 \end{aligned}                                                   \tag{C2.38}
\]
If both restrictions are nontrivial, their $H^0$ and $H^1$
vanish, so $H^1(X_0,\mathcal M)$ is the node fibre.
Frobenius from $\mathcal M$ to $\mathcal M^2$ is a nonzero
semilinear map there. If exactly one restriction is trivial,
its constant section maps isomorphically onto the node fibre,
and $H^1(X_0,\mathcal M)$ is the ordinary elliptic $H^1$
of that component. Frobenius is again bijective.
If both are trivial, the tree condition makes the line
trivial, and $H^1(\mathcal O_{X_0})$ is the direct sum of
the two ordinary elliptic spaces.

For every nontrivial line, $h^0=0$ on both fibres and
$h^1=1$ by Riemann--Roch. Cohomology and base change make
$R^1f_*\mathcal M$ and $R^1f_*(\mathcal M^2)$ free
$R$-modules of rank one. The Frobenius-linearized map
\[
 F_R^*(R^1f_*\mathcal M)\longrightarrow R^1f_*(\mathcal M^2)
\]
has nonzero determinant on the special fibre, hence a unit
determinant over $R$. It is an isomorphism on the geometric
generic fibre. The trivial line gives the same conclusion
in rank two. This handles all $L^4$ arrows simultaneously.
The smooth generic curve is ordinary and every nontrivial
$L$-torsion character arrow is bijective; its theta avoids
$J[L]$. No theta line bundle on the singular curve is assumed.

This finite-level calculation agrees with the opened
generic ordinariness statement
\citep[Proposition 1.3]{A_Yang2016}; its composite-order
primitive-character formulation elsewhere is not used.
Denote the nonempty open stack just obtained by $\mathcal U$.
It is geometrically irreducible. Relatively, theta lies on
the Jacobian of the Frobenius twist; relative Frobenius
identifies all odd torsion Galois-equivariantly, so this
does not change any following level or Frobenius condition.

\proofparagraph{Mixed symplectic level and an actual parameter scheme.}
Put $N=L^2$ and take a finite field $k_0$ containing $\mu_N$.
The complete Jacobi-level theorem
\citep[Theorem 5.15, pp.108--109]{A_DeligneMumford1969}
makes the geometric fibres of the symplectic level-$N$
moduli scheme connected over the chosen cyclotomic base.
The level is invertible and $N\ge3$; the scheme is smooth,
so those fibres are irreducible. Its restriction
$\mathcal U[N]$ remains geometrically irreducible.
This is connectedness for the full composite level,
not a claimed independence of separate $\ell$-adic images.

Let $T\to\mathcal U$ be the scheme of tricanonical frames,
that is, bases of $H^0(C,\omega_C^3)$ of rank five.
It has a universal curve and geometrically connected
$\operatorname{GL}_5$ fibres
\citep[\S5.1 and Lemma 5.2]{A_AchterPries2008}.
Let $Y\to T$ be its scheme of symplectic level-$N$ bases.
This is a genuine finite \'etale torsor under
\[
                             G=\operatorname{Sp}_4(\mathbf Z/N).
\]
The map $Y\to\mathcal U[N]$ has the same smooth connected
frame fibres, so $Y$ is geometrically irreducible.
All schemes and actions descend to a finite extension
of $k_0$. The tricanonical Hilbert construction and its
finite cover give quasi-projective schemes, permitting
effective finite Galois descent. Using frames here avoids
treating the coarse moduli map as a free torsor in the
presence of hyperelliptic automorphisms.

\proofparagraph{The prescribed first Frobenius congruence.}
For every $\ell\in S$, let $u_\ell=1+\ell$ modulo $\ell^2$ and set
\[
 g_\ell=\operatorname{diag}
    (u_\ell,u_\ell,u_\ell^{-1},u_\ell^{-1})
       \in\operatorname{Sp}_4(\mathbf Z/\ell^2).
\]
It has order $\ell$ and
\[
 g_\ell\equiv1\pmod\ell,\qquad
 \frac{g_\ell-1}{\ell}\equiv
               \operatorname{diag}(1,1,-1,-1)\pmod\ell .
\]
The derivative matrix is invertible. The Chinese remainder
theorem gives $g\in G$ of order $L$ with these components.

Twist $Y/k_0$ by the descent cocycle sending arithmetic
Frobenius to right multiplication by $g^{-1}$.
It becomes trivial over the extension of degree $L$,
so finite Galois descent gives a quasi-projective twist
$Y^{(g)}$, geometrically isomorphic to $Y$ and irreducible.
Lang--Weil gives a point on it over $k_0^m$ for every
sufficiently large $m$
\citep[Theorem 7.7.1(ii)--(iii)]{A_Poonen2017}.
Choose $m\equiv1\pmod L$. Under the geometric identification
with $Y$, its level frame $y$ satisfies
\[
                    \operatorname{Frob}_{k_0}^m(y)=y\cdot g .
\]
Its image $t\in T(k_0^m)$ defines an actual smooth curve
$C/k_0^m$, and Frobenius on $J_C[N]$ has matrix $g$
in this frame. The inverse convention gives $g^{-1}$
with the same conclusions. The Weil multiplier is one
because $\mu_N\subset k_0$.

Thus $\pi-1=\ell A_\ell$ with $A_\ell$ an integral
automorphism on every $T_\ell J_C$: its reduction is the
displayed invertible derivative. Since $C\in\mathcal U$,
theta avoids $J[L]$. Equations (C2.36)--(C2.37) prove
(C2.35). Only Lang--Weil on a fixed geometrically irreducible
scheme has been used; no Hilbertian assertion about
a finite field is involved.

\proofparagraph{The covers abelian over the original curve.}
For a geometrically connected \'etale cover $D\to C$
abelian over $C$ with $S$-supported group, every nontrivial
character line has nontrivial $S$-supported order and avoids
theta. Its one-dimensional $H^1$ has a bijective Frobenius
arrow to the doubled character. The invariant part is
ordinary, so $D$ is ordinary. All nontrivial characters
survive, up to duality, in geometric two-torsion.
They separate the abelian deck group, proving faithfulness
already on $J_D[2](\overline\F_2)$.
Every quotient of $J_D$ is ordinary; a nonconstant map to
a two-rank-zero abelian variety would generate a
positive-dimensional ordinary quotient inside it,
which is impossible.

The independent mathematical reviews are
\path{audit/A_finite_prime_theta_root_review.md} and
\path{audit/A_finite_prime_theta_topos_review.md}.
The complete expanded argument and exact source loci are
retained in \path{audit/A_finite_prime_theta_candidate.md}.
This construction permits $C$ to vary with $S$. It does
not assert a curve avoiding theta torsion at all odd
orders, exclude adaptive choices of primes or nonabelian
covers, or produce colliding point sections.
\end{proof}

\subsection{A two-dimensional ordinary genus-two locus detected by degree seven}

\begin{claim}{A.C2.15}{Proved}
Over an algebraically closed field $\Omega$ of characteristic two, put
\[
 u=x(x+1),\qquad q_t=\frac{(x+t)^7}{u},\qquad
 H_t=\frac{u^3}{(x+t)^5(x^2+t)^2},\qquad t\ne0,1.
\]
Choose $r\ne s$ with
\[
 r,s\notin\{0,1,t,\sqrt t\},\qquad H_t(r)=H_t(s).
\]
Define
\[
 \begin{aligned}
 e&=r+s,&\kappa&=e^2H_t(r),&
 a&=\kappa,\\
 b&=\frac{\kappa t^7}{r^2s^2},&
 c&=\frac{\kappa(t+1)^7}{(r+1)^2(s+1)^2}.
 \end{aligned}                                                   \tag{C2.39}
\]
Then the ordinary genus-two curve
\[
               C:\quad y^2+y=\phi(x)
                    =ax+\frac b x+\frac c{x+1}                  \tag{C2.40}
\]
has an explicit theta point of exact order seven. Its associated
cyclic \'etale degree-seven cover has genus eight and two-rank two,
with trivial deck action on geometric two-primary torsion.
Thus A1 holds for its finitely generated forms over $\F_2$,
over every finite extension, and for its hyperbolic open subcurves.

The ordered parameter space is smooth and finite \'etale of degree
$72$ over $(\mathbf A^1\setminus\{0,1\})\times\mathbf G_m$ via
$(t,r,s)\mapsto(t,H_t(r))$. Its coefficient and geometric moduli
maps have finite fibres, and its moduli image has dimension two.
A component has a nonempty open part outside
$(a+b+c)^3+abc=0$, the earlier cubic positive locus.
The construction describes the separable-quadratic-$B$ chart
of ordinary theta-seven incidence.

There is also the nonisotrivial one-parameter family
\[
 U=t(t+1),\qquad
 a=U^{-2},\quad b=\frac{t^3}{(t+1)^2},\quad
 c=\frac{(t+1)^3}{t^2},\qquad t\ne0,1,                           \tag{C2.41}
\]
with the same A1 conclusion. It includes A.C2.13.
\end{claim}
\begin{proof}
\proofparagraph{The parameter cover.}
The degree of $H_t$ is nine: its zeros $0,1,\infty$ have
order three, and its poles $t,\sqrt t$ have orders five and four.
These points are distinct. Differentiation gives
\[
 q_t'=\frac{(x+t)^6(x^2+t)}{u^2},\qquad
 H_t=\frac{q_t}{(q_t')^2},\qquad
 H_t'=\frac{u^2}{(x+t)^6(x^2+t)}.                               \tag{C2.42}
\]
Over
$R=\F_2[t,t^{-1},(t+1)^{-1},\lambda,\lambda^{-1}]$,
the equation $H_t(x)=\lambda$ becomes
\[
                  \lambda(x+t)^5(x^2+t)^2+u^3=0.                \tag{C2.43}
\]
It has degree nine and unit leading coefficient.
None of $x,x+1,x+t,x^2+t$ vanishes in a geometric fibre,
by direct substitution, so all are units in its finite algebra.
Modulo (C2.43), its derivative is
$u^2(x^2+t)/(x+t)$, a unit. Thus it is finite \'etale of
degree nine over the full stated base. The off-diagonal part
of its fibre square is finite \'etale of degree $9\cdot8$.
It is nonempty, smooth and pure of dimension two; no
connectedness assumption about this parameter cover is needed.

\proofparagraph{Principal parts and the universal norm.}
Put
\[
 \begin{aligned}
 \rho_r&=\frac{(r+t)r(r+1)}{r^2+t},&
 \rho_s&=\frac{(s+t)s(s+1)}{s^2+t},&
 B(x)&=(x+r)(x+s),\\
 v(x)&=\frac{\rho_r}{x+r}+\frac{\rho_s}{x+s}
       +\frac{\rho_r}{t+r}+\frac{\rho_s}{t+s},&
 A(x)&=B(x)v(x).
 \end{aligned}
\]
Then $A$ has degree at most two, $B$ is monic quadratic,
and $A(t)=0$. Consider
\[
                       \Phi=\frac{\kappa q_t}{B^2}+v^2+v.
\]
At $r$, use $z=x+r$, so $B=z(e+z)$ and
$B^{-2}=e^{-2}z^{-2}+$ a regular series.
The double and simple principal coefficients of
$\kappa q_t/B^2$ are
$\kappa q_t(r)/e^2$ and $\kappa q_t'(r)/e^2$.
Since $\kappa/e^2=H_t(r)$, (C2.42) makes them
$\rho_r^2,\rho_r$. They cancel exactly against
$(\rho_r/z)^2+\rho_r/z$.
At $s$ the same argument applies because $H_t(s)=H_t(r)$.
There are no cross terms in the square.

Only simple poles at $0,1,\infty$ remain.
Their coefficients are precisely $b,c,a$ of (C2.39);
at infinity $q_t/B^2$ has leading term $x$ and $v$ is
regular. Hence $\Phi-\phi$ is constant. Since
$q_t(t)=v(t)=0$, this constant is $\phi(t)$, proving
\[
                    \frac{\kappa q_t}{B^2}+v^2+v
                               =\phi(x)+\phi(t).                \tag{C2.44}
\]
This principal-part proof applies to every permitted parameter.

Choose $\beta^2+\beta=\phi(t)$ and let
$R_t=(t,\beta)$ and $h=A+B(y+\beta)$. Its conjugate is
$h+B$, so
\[
       \operatorname{Nm}_{C/\mathbf P^1}(h)
          =A^2+AB+B^2(\phi+\phi(t))
          =\frac{\kappa(x+t)^7}{x(x+1)}.                         \tag{C2.45}
\]
The coefficients $a,b,c$ are nonzero. Thus the three
Artin--Schreier poles are simple and have different
exponent two. The curve has genus two, and, as in A.C2.13,
\[
 \operatorname{div}(dx)=2W_0+2W_1-2W_\infty,\quad
 B_C=W_0+W_1-W_\infty,\quad
 \mathcal B_C=\mathcal O(B_C).
\]
The two regular independent logarithmic differentials
$dx/x,dx/(x+1)$ are Cartier-fixed, so the curve is ordinary.
Its Raynaud theta is the symmetric translate
$\{\mathcal O(P-B_C):P\in C(\Omega)\}$.

At $R_t$, $A(t)=0$ and $B(t)\ne0$, so $h$ vanishes
and its conjugate is nonzero. The $x$-map is \'etale there.
Equation (C2.45) gives its sole zero, of exact order seven.
At $W_0,W_1$, the nonzero values of $B$ give simple poles.
At infinity the monic term $x^2y$ has unique largest
pole order five. Consequently
\[
 \operatorname{div}(h)=7R_t-W_0-W_1-5W_\infty,\qquad
 \operatorname{div}(h/u^3)=7(R_t-B_C).                           \tag{C2.46}
\]
The theta point $\eta=\mathcal O(R_t-B_C)$ is nonzero:
otherwise $2R_t$ would be canonical and $R_t$ would be
a hyperelliptic branch point, contrary to $t\ne0,1,\infty$.
It has exact order seven. The normalization of
$z^7=h/u^3$ is therefore geometrically connected and
\'etale, of genus eight.

\proofparagraph{Character cycles and point separation.}
The doubling orbits modulo seven are $\{1,2,4\}$ and
$\{3,6,5\}$. They meet theta at $\eta$ and $-\eta$,
respectively. Since all nontrivial $H^1(C,\eta^i)$ have
dimension one, both Frobenius cycles are nilpotent.
The invariant ordinary part alone survives, of rank two.
Thus the degree-seven cover has two-rank two and its
whole deck group acts trivially on geometric two-torsion,
hence on all two-primary torsion by A.C2.7.

The same claim gives A1: a lifted collision would yield
a two-primary point difference fixed by a nonidentity
free deck transformation, hence distinct linearly
equivalent effective degree-two divisors upstairs.
That would make the cover hyperelliptic, impossible
for a curve admitting a free order-seven automorphism.
Indeed its descended line automorphism has a fixed
point, and an odd-order action on a fibre of at most
two points has a fixed point. All data descend after
a finite extension of the field of definition.
The argument works for finitely generated forms and
every finite extension; the proper-to-open reduction
gives the stated open cases.

\proofparagraph{The image dimension and the converse chart.}
Fix a normalized triple $(a,b,c)$. Every parameter
gives a point $R_t$ with $\mathcal O(R_t-B_C)\in J_C[7]$.
The Abel map is injective, giving at most $7^4-1$
nonzero possibilities. For fixed $R_t$ the divisor
in (C2.46) fixes $h$ up to scalar. The six functions
\[
                         1,x,x^2,y,xy,x^2y
\]
are a basis of $L(W_0+W_1+5W_\infty)$: they are
independent, have the allowed poles, and Riemann--Roch
gives dimension six. The order-five pole makes the
$x^2y$ coefficient nonzero. Requiring $B$ monic
fixes the scalar, and its two ordered roots leave at
most two choices of $(r,s)$. Thus the coefficient
map has at most $2(7^4-1)$ geometric points per fibre.
Every parameter component has two-dimensional image.

An isomorphism between normalized curves preserves
the unique hyperelliptic pencil and permutes
$\{0,1,\infty\}$, allowing only six base transformations.
For each, two simple-pole Artin--Schreier functions
can differ by $w^2+w$ only if $w$ is constant:
any pole of $w$ would give a pole of even order at
least two. A constant changes none of the normalized
coefficients. There are therefore finitely many
coefficient triples per geometric moduli point, so
the moduli image has dimension two.

Conversely, any nonzero theta-seven point on a
normalized ordinary curve gives $h$ with (C2.46).
The basis above makes $B$ monic quadratic; the sole
zero and the poles imply $B(t)B(0)B(1)\ne0$.
If its roots $r,s$ are distinct, the double and simple
principal-part equations in (C2.44) force
$\kappa/e^2=H_t(r)=H_t(s)$ and the displayed residues.
A root with $r^2=t$ would give $q_t'(r)=0$ but
$q_t(r)\ne0$, demanding a zero residue with nonzero
square, which is impossible. The normalization
$v(t)=0$ then recovers this exact chart.
No exhaustion of the double-root boundary is asserted.

For a point outside the cubic locus, let
$\omega^4+\omega+1=0$ and take
\[
             t=\omega,\quad r=\omega^3+\omega,\quad
             s=\omega^3+\omega^2+1 .
\]
Here $H_t(r)=H_t(s)=1$ and
\[
 \begin{aligned}
 (a,b,c)&=(\omega^2+\omega,\,
          \omega^3+\omega^2+1,\,\omega^3+\omega+1),\\
 A&=(\omega^2+\omega)x^2+x+\omega^3+1,\\
 B&=x^2+(\omega^2+\omega+1)x+\omega^3+\omega+1,\\
 (a+b+c)^3+abc&=\omega^2+\omega+1\ne0 .
 \end{aligned}
\]
The component through this point has a nonempty open
part outside the cubic locus; its moduli image still
has dimension two.

\proofparagraph{The separate one-parameter family.}
For (C2.41), choose
$\beta^2+\beta=a(t+1)$ and set
\[
                 h=(x+t)+(x^2+t)(y+\beta).
\]
Put $z=x+t$, so $u=z^2+z+U$.
The identities $b=at^5$, $c=a(t+1)^5$ and $aU^2=1$
give
\[
       \phi(t)=a(t+1),\qquad u(\phi(x)+\phi(t))=z+az^3.
\]
With $A=z$ and $B=z^2+U$,
\[
 \begin{aligned}
 u(A^2+AB)&=z^5+z^3+U^2z,\\
 B^2(z+az^3)&=z^5+az^7+U^2z+z^3 .
 \end{aligned}
\]
Their sum is $az^7$, proving
$\operatorname{Nm}(h)=a(x+t)^7/u$.
The same divisor and character arguments prove A1.
The coefficient $a=U^{-2}$ varies, and the preceding
finite moduli-fibre argument proves nonisotriviality.
Also
\[
                   (a+b+c)^3+abc=\frac{U^3+U+1}{U^3},
\]
so only finitely many parameters lie on the cubic locus.
At $t=\alpha^2$, $\alpha^2+\alpha+1=0$, we get
$(a,b,c)=(1,\alpha,\alpha^2)$, exactly A.C2.13.

The full derivation is retained in
\path{audit/A_theta_seven_surface_candidate.md}.
The original exact verifier is
\path{code/BC_theta_seven_surface_checks.py};
the independent
\path{code/A_theta_seven_surface_independent.py}
uses universal rational identities and a different
quadratic-tower field implementation.
Both pass, with their adjacent JSON files.
The complete reviews are
\path{audit/A_theta_seven_surface_root_review.md} and
\path{audit/A_theta_seven_surface_independent_review.md}.
This proves an explicit positive locus, not existence
of theta seven-torsion on every ordinary genus-two curve.
\end{proof}

\subsection{The complete ordinary theta-five locus}

\begin{claim}{A.C2.16}{Proved}
Over an algebraically closed field of characteristic two, let
\[
 \mathscr P=\{(t,r)\in\mathbf A^2:
          t(t+1)r(r+1)(r+t)(r^2+t)\ne0\}.
\]
For $(t,r)\in\mathscr P$ put
\[
 \begin{aligned}
 d&=\frac{r(r+1)}{r^2+t},&
 \kappa&=\frac{r^3(r+1)^3}{(r+t)^3(r^2+t)^2},\\
 a&=\kappa,&b&=\frac{\kappa t^5}{r^2},&
 c&=\frac{\kappa(t+1)^5}{(r+1)^2}.
 \end{aligned}                                                   \tag{C2.47}
\]
The ordinary genus-two curve
\[
               C_{t,r}:\quad y^2+y=\phi(x)
                     =ax+\frac b x+\frac c{x+1}                 \tag{C2.48}
\]
has an explicit theta point of exact order five.
Its cyclic \'etale degree-five cover has genus six and
two-rank two, with trivial geometric deck action on
two-primary torsion. Hence A1 holds for every finitely
generated form over $\F_2$, after every finite extension,
and for every hyperbolic open subcurve.

Conversely, a normalized ordinary curve (C2.48), with
$abc\ne0$, has nonzero theta five-torsion if and only if
its coefficients arise from (C2.47).
The irreducible parameter surface has finite geometric
fibres over coefficient space and over geometric moduli,
and its moduli image has dimension two. A nonempty open
part lies outside $(a+b+c)^3+abc=0$.
The explicit example below has geometric automorphism
group exactly the order-two hyperelliptic group.
\end{claim}
\begin{proof}
All three coefficients are nonzero. The three simple
Artin--Schreier poles at $0,1,\infty$ have different
exponent two, giving genus two. As in A.C2.13, the
logarithmic differentials $dx/x,dx/(x+1)$ are a regular
Cartier-fixed basis, so the curve is ordinary, and
\[
 \mathcal B_C=\mathcal O(B_\theta),\qquad
 B_\theta=W_0+W_1-W_\infty,\qquad
 \Theta=\{\mathcal O(R-B_\theta):R\in C(\Omega)\}.
\]
Twice $B_\theta$ is canonical. The same locally exact
differential convention and Frobenius twist as in
A.C2.5 apply; odd torsion transports canonically.

\proofparagraph{The norm and its divisor.}
Set $u=x(x+1)$, $\delta=\phi(t)$ and choose
$\beta^2+\beta=\delta$. Put
\[
 R=(t,\beta),\qquad A=d(x+t),\qquad B=x+r,\qquad
                          h=A+B(y+\beta).
\]
The universal polynomial identity is
\[
          u(A^2+AB)+B^2(u\phi+\delta u)
                           =\kappa(x+t)^5.                     \tag{C2.49}
\]
To derive it, write
\[
 q_t=\frac{(x+t)^5}{u},\qquad
 q_t'=\frac{(x+t)^4(x^2+t)}{u^2},\qquad
 \rho=\frac{q_t(r)}{q_t'(r)}
      =\frac{(r+t)r(r+1)}{r^2+t}.
\]
Both $q_t(r)$ and $q_t'(r)$ are nonzero by the parameter
exclusions. Let
$v=\rho/(x+r)+\rho/(t+r)$.
Then $Bv=A$ and $\kappa=q_t(r)/q_t'(r)^2$.
At $r$, the double and simple coefficients of
$\kappa q_t/B^2$ are $\rho^2,\rho$.
They cancel against $v^2+v$. The remaining simple
poles have coefficients $b,c,a$ at $0,1,\infty$,
respectively. The expression vanishes at $t$,
so it equals $\phi(x)+\phi(t)$, proving (C2.49).

The Artin--Schreier conjugate of $h$ is $h+B$.
Thus
\[
             \operatorname{Nm}(h)=\frac{\kappa(x+t)^5}{u}.
\]
At $R$, $h=0$ and its conjugate is $t+r\ne0$.
The $x$-map is \'etale there, so its sole zero has
order five at $R$. The nonzero values $B(0)=r$ and
$B(1)=r+1$ give simple poles at $W_0,W_1$.
At infinity $xy$ has unique largest pole order three;
the remaining terms have pole order at most two.
Consequently
\[
 \operatorname{div}(h)=5R-W_0-W_1-3W_\infty,\qquad
 \operatorname{div}(h/u^2)=5(R-B_\theta).                        \tag{C2.50}
\]
The theta point $\eta=\mathcal O(R-B_\theta)$ is
nonzero: otherwise $2R$ would be canonical, making
$R$ Weierstrass despite $t\notin\{0,1,\infty\}$.
It has exact order five.

\proofparagraph{The cover and the full section consequence.}
The normalization of $z^5=h/u^2$ is geometrically
connected by exact order, and \'etale because every
valuation is divisible by five. Its genus is six.
The nonzero characters form the single doubling
orbit $1,2,4,3$ modulo five. Each $H^1(C,\eta^i)$
has dimension one, and $\eta\in\Theta$ makes one
Frobenius arrow zero. The entire new cycle is
nilpotent. Only the ordinary invariant rank two
survives, so the cover has two-rank two and its
full order-five deck group acts trivially on
geometric two-primary torsion.

Claim A.C2.7 excludes a collision: after lifting,
a nonidentity deck element would fix a two-primary
difference between points over distinct base points.
This gives distinct linearly equivalent effective
degree-two divisors on the upstairs curve.
That curve is not hyperelliptic, since a free
odd-order automorphism would descend to the line
and fix a point in a fibre of size at most two
above a fixed base coordinate. The equivalence
is therefore impossible.

For a finitely generated form, all parameters,
the isomorphism, cover and deck action descend
after a finite extension of the field of definition.
A section collision restricts to that extension;
after defining a chosen lift, A.3 supplies a
colliding lift with two-primary difference.
The same contradiction applies. The proper-to-open
comparison gives the hyperbolic open case.
This retains the full geometrically prime-to-two
groups, not just their abelianizations.

\proofparagraph{The converse has no omitted chart.}
Suppose a normalized ordinary curve with $abc\ne0$
has a theta point $\eta=\mathcal O(R-B_\theta)$
of exact order five. A Weierstrass point $W$ has
$2W\sim2B_\theta$, so its theta class is two-torsion;
thus $R=(t,\beta)$ with $t\ne0,1$ and
$\beta^2+\beta=\phi(t)$.
Choose $g$ with divisor $5(R-B_\theta)$ and
put $h=gu^2$. It has exactly the first divisor
in (C2.50). Riemann--Roch gives
\[
             L(W_0+W_1+3W_\infty)=\langle1,x,y,xy\rangle.
\]
The divisor has degree five and its space dimension
four; these four independent functions obey the
pole bounds. Hence $h=A+B(y+\beta)$ with $A,B$
linear. The order-three infinity pole forces
$B$ to have degree one. Scale it to $B=x+r$.
Since $h(R)=0$, write $A=d(x+t)$.

The poles force $r\ne0,1$. Also $r\ne t$:
otherwise both $A$ and $B$ vanish at $t$, making
$h$ vanish at both unramified points above $t$.
Its norm is $\kappa q_t$ for a nonzero constant,
and its leading term gives $\kappa=a$.
Evaluation at $r$, followed by differentiating
the norm identity and evaluating there, yields
\[
                    A(r)^2=\kappa q_t(r),\qquad
                    A(r)=\kappa q_t'(r).                       \tag{C2.51}
\]
The first makes $A(r)\ne0$, so the second makes
$q_t'(r)\ne0$ and forces $r^2\ne t$.
Their quotient gives
\[
       \kappa=\frac{q_t(r)}{q_t'(r)^2},\qquad
       d=\frac{\kappa q_t'(r)}{r+t}
                        =\frac{r(r+1)}{r^2+t}.
\]
Residues at $0,1$ give exactly $b,c$ in (C2.47).
Thus every nonzero ordinary theta-five point
belongs to the asserted parameter space.

\proofparagraph{Dimension and a curve with no odd automorphisms.}
The space $\mathscr P$ is a nonempty irreducible
open subset of $\mathbf A^2$.
For a fixed normalized coefficient triple, the
injective Abel map confines $R$ to the finite set
$J[5]$. Its divisor fixes $h$ up to scalar,
and making its linear $y$-coefficient monic
determines $h$ and $r$ uniquely. The coefficient
map therefore has finite geometric fibres and
two-dimensional image.
As in A.C2.15, an isomorphism must permute the
three hyperelliptic branch coordinates, giving
only six base maps. An Artin--Schreier correction
preserving only simple poles must be constant.
There are consequently finitely many normalized
triples per geometric isomorphism class, and
the moduli image has dimension two.

Let $\omega^4+\omega+1=0$ and take
$t=\omega$, $r=\omega^2$.
Then
\[
 \begin{aligned}
 (a,b,c)&=(1,\omega,\omega+1),&
 d&=\omega^2+\omega+1,\\
 \delta&=\omega,&\beta&=\omega^3+\omega,&
 (a+b+c)^3+abc&=\omega^2+\omega\ne0 .
 \end{aligned}
\]
In particular a nonempty open part of this irreducible
surface lies outside the cubic theta locus.
All three coefficients of this example are distinct.
The generating branch changes act, up to a removable
constant, by
\[
 x\mapsto x+1:\ (a,b,c)\mapsto(a,c,b),\qquad
 x\mapsto1/x:\ (a,b,c)\mapsto(b,a,c).
\]
No nontrivial permutation fixes this triple.
An automorphism over the identity base map is an
automorphism of the separable quadratic extension,
hence either the identity or $y\mapsto y+1$.
More generally, a possible rescaling in an
Artin--Schreier generator is forced to be one:
the coefficient of $y$ requires $e(x)^2+e(x)=0$.
Thus the geometric automorphism group is exactly
the hyperelliptic group of order two.
This example has no nontrivial odd-order
geometric automorphism and no odd-order Galois
elliptic quotient. Semilinear constant-field
automorphisms are not included in this assertion.

The universal identities in $\F_2(x,t,r)$ and
the exact finite example are checked by
\path{code/A_theta_five_surface_checks.py}, with
adjacent JSON. All $168$ allowed pairs over
$\F_{16}$ pass, of which $156$ lie outside the
cubic locus; the family proof uses the universal
identity, not these finite counts.
The complete derivation and independent review are
\path{audit/A_theta_five_surface_candidate.md} and
\path{audit/A_theta_five_surface_degeneration_review.md}.
No existence of theta five-torsion on every
ordinary genus-two curve, or universal A1,
is inferred.
\end{proof}




\subsection{A primitive Prym can suffice when the deck action is faithful}

\begin{claim}{A.C2.17}{Proved}
Let $K$ be finitely generated over $\F_2$, and let $C/K$ be a
smooth proper geometrically connected curve of genus two, with
Raynaud theta curve $\Theta$ as in A.C2.5. Suppose that over
$\overline K$ there is a line bundle $L$ of exact order
\[
 n=\ell^a m^b,\qquad
 \ell\ne m\text{ odd primes},\quad a,b\ge1,
\]
such that every multiplication-by-two orbit on
$(\Z/n)^\times$ contains an index $i$ with $L^i\in\Theta$.
Then the geometrically prime-to-two point-section map of $C$
is injective. No condition is needed on characters of smaller
conductor. The conclusion holds after every finite extension,
for forms satisfying this geometric hypothesis, and for every
hyperbolic open subcurve of $C$.

In particular, an exact-order-fifteen theta point suffices,
without any assumption on theta three- or five-torsion.
Every exact theta point of order $3^a5^b$, with $a,b\ge1$,
also suffices. Combining this with A.C2.8, for any odd $n>1$
the simple sufficient condition is
\[
 L\in\Theta\text{ of exact order }n,\qquad
             \langle2,-1\rangle=(\Z/n)^\times.             \tag{PP.1}
\]
These assertions do not assume that $C$ is ordinary.
\end{claim}
\begin{proof}
We first prove a geometric lemma. Let $k$ be algebraically
closed of characteristic two, and let $D/k$ have genus
$n+1$, where $n\ge15$ is odd. Suppose that a cyclic group
$G$ of order $n$ acts freely on $D$ and that
\[
                            f(D)<n-1.                     \tag{PP.2}
\]
Then $D$ has no nonconstant rational function of degree at
most four.

We use the Castelnuovo--Severi inequality in the following
exact form. If $F=F_1F_2$ is a compositum of function
fields of curves over a perfect field, $d_i=[F:F_i]$, and
their genera are $g,g_1,g_2$, then
\[
              g\le d_1g_1+d_2g_2+(d_1-1)(d_2-1).
                                                               \tag{PP.3}
\]
The opened primary author version of
\citet[Proposition 2.1, author p.~3]{A_PrimPoonen2007}
states this without a restriction on the characteristic.
All composita and degrees in its applications are
specified below.

The curve $D$ is not hyperelliptic. If it were, its
hyperelliptic involution would be unique; uniqueness also
follows from (PP.3), since distinct index-two rational
subfields would give $g(D)\le1$. Hence $G$ descends to
its genus-zero quotient. The action there is faithful,
as its kernel lies in the order-two hyperelliptic group.
A generator of an odd cyclic subgroup of
$\operatorname{PGL}_2(k)$ has a fixed point. Its fibre
on $D$ has degree two and is $G$-invariant. Freeness
partitions its geometric points into orbits of size $n$,
with constant multiplicity within each orbit, making
its degree divisible by $n$. This is impossible.

More generally, a $G$-invariant rational subfield
$k(z)\subset k(D)$ of degree $d<n$ is impossible. If
the induced cyclic action on $\mathbf P^1$ is trivial,
every fibre is invariant. If it is nontrivial, a fibre
above a fixed point is invariant. In either case
freeness would force $n\mid d$.

Degree one is excluded by the genus, and degree two by
nonhyperellipticity. A degree-three pencil is unique:
two distinct index-three rational subfields generate
$k(D)$ and would give $g(D)\le4$ by (PP.3).
Uniqueness makes this pencil $G$-invariant, again
impossible. There is no inseparability exception:
over $k$, an inseparable rational function is a square.
If its degree is at most four, a square root has degree
at most two, which has just been excluded.

Suppose therefore that $z$ has degree four. It is
separable. Its rational subfield cannot be
$G$-invariant, so for some $\delta\in G$ the fields
$k(z)$ and $k(\delta z)$ differ. Put
\[
                 F=k(z,\delta z),\qquad e=[k(D):F].
\]
Then $e\mid4$. The value $e=4$ would make the two
rational subfields equal. If $e=1$, (PP.3) gives
$g(D)\le9$, contrary to $g(D)=n+1\ge16$.
Thus $e=2$. The smooth curve $E$ with function field
$F$ has two degree-two rational maps generating its
function field, so (PP.3) gives $g(E)\le1$.
Genus zero would make $D$ hyperelliptic. Consequently
$g(E)=1$. The extension $k(D)/F$ is separable,
because $k(D)/k(z)$ is separable, and has degree two.
Write $\tau$ for its involution.

This bielliptic involution is unique. Two distinct
degree-two elliptic subfields generate $k(D)$, and
(PP.3) with $d_1=d_2=2$, $g_1=g_2=1$ would give
$g(D)\le5$. Hence every automorphism of $D$ commutes
with $\tau$, in particular every element of $G$.

The fixed-point set of $\tau$ is nonempty:
Riemann--Hurwitz gives total different $2n$ for
$D\to E$. It is $G$-stable, and $G$ acts freely on it.
Its cardinal $r$ is a positive multiple of $n$.
Every ramification point of a separable quadratic
extension in characteristic two has different exponent
at least two, since its zeroth and first lower
ramification groups both have order two.
It follows that $2r\le2n$, so $r=n$.
The ramified Deuring--Shafarevich formula now gives
\[
 f(D)-1=2(f(E)-1)+n,\qquad
                   f(D)=2f(E)+n-1\ge n-1,                 \tag{PP.4}
\]
contradicting (PP.2). The exact source is
\citet[Corollary 1.8, pp.~35--36]{A_PrimCrew1984};
both the ramified statement and its proof were opened.

There is also an elementary proof of the needed lower
bound, which makes its geometric content explicit.
The action of $G$ on $E$ is faithful because its
kernel is contained in $\langle\tau\rangle$.
A nonidentity element of $G$ cannot fix a point on
$E$: its odd-order permutation of an at-most-two-point
fibre would be trivial, giving a fixed point on $D$.
Every fixed-point-free elliptic automorphism is a
translation. Indeed, after choosing an origin write
it as $T_b\circ u$ with $u$ fixing the origin. If
$u\ne1$, the nonzero endomorphism $1-u$ is surjective
and produces a fixed point. Thus $G$ acts by an
order-$n$ group of translations.

Take one ramification orbit $P_0,\ldots,P_{n-1}$
and let $b_i$ be its images on $E$.
The differences $b_i-b_0$ are killed by $n$.
Pullback of divisors, whose ramification indices are
two, shows that
\[
                       T_i=n[P_i-P_0],\quad 1\le i<n,
\]
are killed by two. They are linearly independent over
$\F_2$. Otherwise a nonempty relation with
$\epsilon_i\in\{0,1\}$ would give
\[
 \operatorname{div}(v)=
 n\left(\sum_{i=1}^{n-1}\epsilon_iP_i
       -\left(\sum_{i=1}^{n-1}\epsilon_i\right)P_0\right).
\]
Its divisor is $\tau$-invariant, so $\tau(v)/v$ is
constant. The constant has square one and therefore
equals one in characteristic two. Hence $v\in k(E)$
and every valuation at $P_i$ is even. A nonzero
$\epsilon_i$ gives valuation $n$, which is odd.
This contradiction proves $f(D)\ge n-1$
independently of the formula. The geometric lemma is
proved.

Return to the cyclic \'etale cover $D\to C$ associated
with $L$, and let $\sigma$ generate its deck group.
Its genus is $n+1$. Each nonzero character space
$H^1(C,L^i)$ has dimension one. A.C2.5 shows that
the stated theta intersections make every primitive
doubling orbit Frobenius-nilpotent. There are
$\varphi(n)$ primitive characters, so
\[
 f(D)\le f(C)+n-1-\varphi(n)
       \le n+1-\varphi(n)<n-1,                            \tag{PP.5}
\]
since $\varphi(n)\ge8$. The geometric lemma applies.

Set
\[
 a_0=\sigma^{n/\ell},\qquad b_0=\sigma^{n/m},\qquad
          A=(1-a_0)(1-b_0)\in\operatorname{End}(J_D).
                                                               \tag{PP.6}
\]
These automorphisms have distinct odd prime orders
$\ell,m$. On a character of index $i$, the scalar of
$A$ is nonzero exactly when neither $\ell$ nor $m$
divides $i$, equivalently when $i$ is primitive
modulo $n$. It follows that $A$ kills all geometric
two-primary torsion, with the following integral
justification.

Pass from $\Z_2$ to a finite unramified extension
containing all $n$-th roots of unity. Since $n$ is odd,
the cyclic character idempotents act integrally on
the \'etale Tate module. Their summands are free,
and each rank equals the dimension after reduction
modulo two. The Artin--Schreier character description
in A.C2.5 and A.C2.8 makes every primitive summand
zero. Duality may invert the character indices,
but it preserves primitivity.
On each remaining summand one factor of
(PP.6) is exactly zero. Thus $A$ vanishes on
$T_2J_D$ and on $J_D[2^\infty](\overline K)$.
This is not an inference from an arbitrary
endomorphism being zero modulo two.
Equivalently, $\operatorname{im}(A)$ is an abelian
subvariety of two-rank zero, although the
smaller-conductor factors may be ordinary.

If distinct $P,Q\in C(K)$ gave equal section classes,
make a finite extension defining $L,D,\sigma$ and a
chosen lift $P'$. By A.3 there is a lift $Q'$ with
$[P'-Q']$ two-primary. Applying (PP.6) gives
\[
 P'+a_0b_0P'+a_0Q'+b_0Q'
 \ \sim\
 a_0P'+b_0P'+Q'+a_0b_0Q'.                                \tag{PP.7}
\]
The elements $1,a_0,b_0,a_0b_0$ are distinct, because
the two prime-order subgroups intersect trivially.
The free deck action and the distinct base points
make all eight displayed points distinct.
Thus (PP.7) gives a rational function of degree four,
with simple zeros and poles; in particular it is
separable. This contradicts the geometric lemma.
All geometric arguments are over $\overline K$,
while the lifted two-primary difference is over a
finite extension by A.3. No specialization of the
theta hypothesis is needed. Restriction and
functoriality to the compactification give the
stated forms and open-subcurve assertions.

It remains to check the arithmetic corollaries.
The divisor $\Theta$ is symmetric, so a theta point
$L$ supplies both $L$ and $L^{-1}$. Under (PP.1)
each primitive doubling orbit contains $1$ or $-1$,
which proves the requisite orbit condition.
For $n=15$ the primitive orbits are
\[
              \{1,2,4,8\},\qquad \{14,13,11,7\}.
\]
Hence $f(D)\le8$, whereas a bielliptic alternative
would force $f(D)\ge14$. No smaller-order
theta exclusion is necessary.

For $a,b\ge1$,
\[
 \operatorname{ord}_{3^a}(2)=2\cdot3^{a-1},\qquad
 \operatorname{ord}_{5^b}(2)=4\cdot5^{b-1}.
\]
For the first identity, $v_3(2^2-1)=1$, and expansion
of $(1+3^ru)^3-1$ with $3\nmid u$ raises its valuation
from $r$ to $r+1$. The second follows in the same
way from $v_5(2^4-1)=1$ and
$(1+5^ru)^5-1$. Reduction modulo three and five
gives the exact initial orders. The displayed
orders have greatest common divisor two, so two
has order $\varphi(3^a5^b)/2$ modulo $3^a5^b$.
Moreover $-1$ is not a power of two there:
modulo three its exponent would be odd, whereas
modulo five it would be $2\bmod4$.
Thus $\langle2,-1\rangle$ is the full unit group.

Finally, if an odd $n$ has $r$ distinct prime
divisors, the Chinese remainder theorem gives
\[
  (\Z/n)^\times/\bigl((\Z/n)^\times\bigr)^2
                    \simeq(\Z/2)^r.
\]
Two generators cannot span this quotient when
$r>2$. Condition (PP.1) therefore forces at most
two prime divisors. The one-prime case is A.C2.8
and the two-prime case is the theorem just proved.

At degree fifteen the surviving order-three and
order-five characters may have trivial common
kernel. The endomorphism (PP.6) still kills both
and produces a degree-four divisor relation,
explaining the improvement over A.C2.7.
This does not supply a suitable theta point on
every curve: torsion density and A.C2.14 leave
the varying-order existence question open.
The exact arithmetic certificate
\path{code/BC_two_prime_prym_checks.py} and its
JSON output record 289 passing checks, including
261 two-prime orders and eight exact cyclotomic
identities. The proof above establishes the
unbounded assertions.
\end{proof}


\subsection{An exact mixed-order theta point and all split ordinary models over \texorpdfstring{$\F_4$}{F4}}

\begin{claim}{A.C2.18}{Proved}
Let $\alpha^2+\alpha+1=0$ and
\[
 C/\F_4:\qquad y^2+y=x+\frac{\alpha}{x}
                         +\frac{\alpha}{x+1}.                 \tag{C2.52}
\]
Its Raynaud theta has no points killed by three or five,
but has an explicit point of exact order fifteen.
The resulting cyclic \'etale cover has genus sixteen
and two-rank eight. Its geometric deck action on
two-primary torsion is faithful, but its primitive
conductor-fifteen part has two-rank zero.
Claim A.C2.17 therefore proves A1 for this curve,
after every finite extension, for its finitely generated
forms, and for its hyperbolic open subcurves.

Together with A.C2.1, A.C2.11, A.C2.13 and A.C2.15,
this proves A1 for every ordinary genus-two curve over
$\F_4$ whose three hyperelliptic ramification points
are rational, with the same extension and form scope.
This finite enumeration does not presume rationality
of the branch points for every ordinary $\F_4$ curve.
\end{claim}
\begin{proof}
Let $W_0,W_1,W_\infty$ be the three ramification points,
$u=x(x+1)$ and $B_\theta=W_0+W_1-W_\infty$.
As in A.C2.13, the different calculation gives
$\operatorname{div}(dx)=2B_\theta$ and identifies the
locally exact differential bundle with
$\mathcal O(B_\theta)$. The two logarithmic differentials
$dx/x,dx/(x+1)$ prove ordinariness.
Thus the symmetric Raynaud theta is
\[
          \Theta=\{\mathcal O(R-B_\theta):R\in C(\overline{\F}_2)\},
          \qquad 0\notin\Theta.
\]
The cubic obstruction of A.C2.2 is
$(1+\alpha+\alpha)^3+\alpha^2=\alpha\ne0$.
Hence $\Theta\cap J[3]=\varnothing$.

\proofparagraph{A complete geometric Hankel criterion.}
For any ordinary normalized curve
$y^2+y=\phi(x)=ax+b/x+c/(x+1)$ with $abc\ne0$,
let $n=2d+3\ge3$ be odd and
$R=(t,\beta)$ with $t\ne0,1$ and $\beta^2+\beta=\phi(t)$.
Then $n[R-B_\theta]=0$ if and only if a nonzero
\[
 h\in L\bigl(W_0+W_1+(n-2)W_\infty\bigr)
\]
vanishes to order at least $n$ at $R$.
The pole bound has degree $n$, so such a function
necessarily has exactly the divisor
$nR-W_0-W_1-(n-2)W_\infty$.
Dividing by $u^{(n-1)/2}$ gives $n(R-B_\theta)$;
multiplication by that power proves the converse.
Riemann--Roch gives the pole-space basis
\[
                 1,x,\ldots,x^d,\ y,xy,\ldots,x^dy.
\]
The $2d+2=n-1$ independent displayed functions satisfy
the bounds, and $n>2g-2$ gives exactly this dimension.
Thus write $h=A+B(y+\beta)$ with $\deg A,\deg B\le d$.
A ramification point supplies no omitted odd theta point:
$2W\sim2B_\theta$ makes its class two-torsion,
and $0\notin\Theta$.

Put $z=x+t$ and $y+\beta=\sum_{i\ge1}y_i z^i$.
The Artin--Schreier equation gives the exact recursion
\[
 y_i=a\delta_{i1}+\frac b{t^{i+1}}+\frac c{(t+1)^{i+1}}
       +\begin{cases}y_{i/2}^2,&2\mid i,\\0,&2\nmid i.\end{cases}
                                                               \tag{C2.53}
\]
For $B(t+z)=\sum_{j=0}^d b_jz^j$, the coefficients
of $A$ uniquely cancel jets through order $d$.
The remaining equations are
\[
             \sum_{j=0}^d b_jy_{i-j}=0,
                         \qquad d+1\le i\le n-1.              \tag{C2.54}
\]
If $B=0$, a polynomial $A$ of degree at most $d$
cannot vanish to order $n$, so no spurious nonzero
solution occurs. With $D=t(t+1)$, set $P_i=D^{2i}y_i$.
These are polynomials:
\[
 P_i=bt^{i-1}(t+1)^{2i}+ct^{2i}(t+1)^{i-1}
       +a\delta_{i1}D^2
       +\begin{cases}P_{i/2}^2,&2\mid i,\\0,&2\nmid i.\end{cases}
                                                               \tag{C2.55}
\]
Multiplying row $i$ of (C2.54) by $D^{2i}$ and
column $j$ by $D^{-2j}$ replaces its entries by
$P_{i-j}$. These are invertible operations away
from $t=0,1$. Therefore the geometric theta-$n$
coordinates are exactly the common zeros there
of all maximal minors of the $(d+2)\times(d+1)$
polynomial matrix $(P_{i-j})$.
Not all minors vanish identically: otherwise the
injective Abel map would place an entire curve
in the finite group $J[n]$.
Taking their polynomial gcd and deleting every
factor $t$ or $t+1$ is consequently a complete
algebraic-closure test, rather than a bounded
rational-point search. A partial gcd already
constant after these deletions proves emptiness;
a positive conclusion requires all maximal minors.

\proofparagraph{The geometric five-torsion exclusion.}
For (C2.52), use a variable $\xi$ for the proposed
coordinate. Formula (C2.55) gives
\[
\begin{aligned}
 P_1&=\xi^4+\xi^2+\alpha,\\
 P_2&=\xi^8+\alpha^2\xi^4+\alpha\xi+\alpha^2,\\
 P_3&=\alpha(\xi^4+\xi^2),\\
 P_4&=\xi^{16}+\alpha\xi^{10}+\alpha\xi^9
                         +\alpha\xi^3+\alpha^2\xi^2+\alpha.
\end{aligned}
\]
The relevant matrix is
$\bigl(\begin{smallmatrix}P_2&P_1\\P_3&P_2\\P_4&P_3\end{smallmatrix}\bigr)$.
Its first minor equals
\[
                 P_2^2+P_1P_3=(\xi^4+\xi+\alpha)^4.
\]
Modulo $f=\xi^4+\xi+\alpha$, writing $D=\xi^2+\xi$,
the four entries reduce to
$(D,D+1,\alpha D+\alpha^2,\alpha D+1)$.
The other minor $P_2P_3+P_1P_4$ reduces to
$\alpha^2\ne0$. Thus the minors have no common
zero, and $\Theta\cap J[5]=\varnothing$.
This excludes points over every algebraic extension.

\proofparagraph{The exact divisor of order fifteen.}
Choose $t^2+\alpha t+1=0$, so that
$t\in\F_{16}$, $t^4+t=\alpha$ and $t\ne0,1$.
Let $\beta^2+\beta=\alpha t$ and $R=(t,\beta)$.
Indeed $\phi(t)=\alpha t$; its absolute trace
from $\F_{16}$ is one, and $\beta$ may be taken
in $\F_{256}$. Define $A,B$ by their coefficient
lists in ascending powers of $x$:
\[
\begin{aligned}
 A={}&[\alpha^2+t,\ \alpha t,\ t,\ 1+t,
           \ \alpha^2+\alpha t,\ 1+\alpha^2t,\ 1+t],\\
 B={}&[\alpha,\ \alpha+t,\ \alpha+t,\ 1+t,
                  \ \alpha^2t,\ \alpha t,\ 1].
\end{aligned}                                                   \tag{C2.56}
\]
With $F=x^3+x^2+\alpha=u\phi(x)$, direct multiplication
in $\F_2[\alpha,t]/(\alpha^2+\alpha+1,t^2+\alpha t+1)$
gives
\[
          u(A^2+AB)+B^2(F+\alpha t\,u)=(x+t)^{15}.              \tag{C2.57}
\]
The same calculation gives
$A(t)=0$, $B(t)=t$, and $B(0)=B(1)=\alpha$.
For $h=A+B(y+\beta)$, its conjugate is $h+B$,
and (C2.57) gives
$\operatorname{Nm}(h)=(x+t)^{15}/u$.
At $R$, $h=0$ and its conjugate is $t\ne0$.
The $x$-map is \'etale there, so $h$ has its
sole zero of order fifteen at $R$.
The nonzero coefficients at $0,1$ give a simple
pole at each of $W_0,W_1$.
At infinity the monic term $x^6y$ has the unique
largest pole order thirteen; $A$ and $\beta B$
have pole order at most twelve. Consequently
\[
 \begin{aligned}
 \operatorname{div}(h)&=15R-W_0-W_1-13W_\infty,\\
 \operatorname{div}(h/u^7)&=15(R-B_\theta).
 \end{aligned}                                                   \tag{C2.58}
\]
The theta class $L=\mathcal O(R-B_\theta)$ has
exact order fifteen, since orders one, three
and five have all been excluded.
The normalization of $z^{15}=h/u^7$ is therefore
geometrically connected and \'etale. After a finite
constant extension it is cyclic Galois; all fifteenth
roots of unity lie in $\F_{16}$.
Its genus is $1+15(2-1)=16$.

\proofparagraph{Faithful deck action and a vanishing primitive part.}
The nonzero doubling orbits modulo fifteen are
\[
 \{1,2,4,8\},\quad\{7,14,13,11\},\quad
                   \{3,6,12,9\},\quad\{5,10\}.                 \tag{C2.59}
\]
The first two are primitive and are negatives
of one another. They meet theta at $L$ and
$L^{-1}$; each contains a zero Frobenius arrow
between one-dimensional character spaces.
Their entire Frobenius blocks are nilpotent.
The last two blocks survive because theta
contains no nonzero five- or three-torsion.
The invariant ordinary rank is two, so
\[
                             f(D)=2+4+2=8.
\]
Surviving characters of indices three and five
have kernels of orders three and five, with
trivial intersection. The geometric deck action
on two-torsion, and hence on its Tate module,
is faithful. Nevertheless
\[
                    (1-\sigma^3)(1-\sigma^5)
\]
kills all geometric two-primary torsion.
For colliding lifts $P',Q'$ over distinct
points of $C$, applying this operator to
$[Q'-P']$ would give the distinct effective
four-divisor equivalence
\[
 Q'+\sigma^8Q'+\sigma^3P'+\sigma^5P'
       \sim\sigma^3Q'+\sigma^5Q'+P'+\sigma^8P'.                 \tag{C2.60}
\]
The divisors are distinct because the two free
deck orbits are disjoint. The geometric lemma
in A.C2.17 excludes every rational function
of degree at most four on this genus-sixteen
curve of two-rank eight. Thus the equivalence
is impossible. Equivalently, apply A.C2.17
directly to the two primitive orbits.
This proves A1 for the stated curve, with
all the extension, form and open-subcurve
consequences specified in that claim.

\proofparagraph{All ten split normalized $\F_4$ triples.}
The branch changes $x\mapsto x+1$ and $x\mapsto1/x$
permute $a,b,c$, up to a removable constant.
There are exactly ten unordered nonzero triples
from $\{1,\alpha,\alpha^2\}$:
\begin{center}
\begin{tabular}{ll}
\toprule
Triples, with indicated conjugates & Separating construction\\
\midrule
$(1,1,1),(\alpha,\alpha,\alpha),(\alpha^2,\alpha^2,\alpha^2)$
  & A.C2.1, cubic\\
$(1,1,\alpha),(1,1,\alpha^2)$ & A.C2.11, ninth order\\
$(1,\alpha,\alpha^2)$ & A.C2.13, seventh order\\
$(\alpha,\alpha,\alpha^2),(\alpha,\alpha^2,\alpha^2)$
  & A.C2.15, seventh order\\
$(1,\alpha,\alpha),(1,\alpha^2,\alpha^2)$
  & this claim, fifteenth order\\
\bottomrule
\end{tabular}
\end{center}
For the fourth row, in the one-parameter family
of A.C2.15 choose $t^2+t=U$ with
$U=\alpha^2$ or $\alpha$. Then
$t^4=t+1$, $t^5=(t+1)^5=U$, and its coefficients
are $a=U^{-2}$, $b=c=U^{-1}$, as required.
Every ordinary genus-two curve over $\F_4$ with
three rational ramification points has a rational
hyperelliptic quotient and a normalization of
this form: its reduced Artin--Schreier poles
are simple and can be placed at $0,1,\infty$.
A constant term is removed after a finite
extension, which is allowed throughout.
No assertion about nonsplit branch triples
is needed for this finite enumeration.

The exact identity, all local values, the
five-torsion exclusion and the character
calculations are checked by
\path{code/A_theta_fifteen_certificate.py}
and, using independent quotient-polynomial
arithmetic and formal jets, by
\path{code/A_theta_fifteen_root_independent.py},
with their adjacent JSON files.
The complete reviews are
\path{audit/A_two_prime_prym_root_review.md} and
\path{audit/A_theta_fifteen_root_review.md}.
The broader bounded scans in
\path{audit/A_theta_Hankel_scan.md} are retained
as discovery evidence. They imply no unbounded
prime-power exclusion or universal torsion
cutoff, and universal A1 remains A.30.
\end{proof}

\subsection{Primitive character vanishing at an arbitrary odd conductor}

\begin{claim}{A.C2.19}{Proved}
Let $K$ be finitely generated over $\F_2$, and let $C/K$
be smooth proper geometrically connected of genus two.
Suppose that, geometrically, a line bundle $L$ of exact
odd order $n>1$ has the following property: every
multiplication-by-two orbit on $(\Z/n)^\times$ meets
\[
                         \{i:L^i\in\Theta\}.
\]
Then the geometrically prime-to-two point-section map
of $C$ is injective. The assertion holds over every
finite extension, for forms satisfying this geometric
hypothesis, and for hyperbolic open subcurves.
There is no restriction on the number of prime
divisors of $n$, and $C$ need not be ordinary.

Equivalently, it suffices that the primitive
conductor-$n$ factor of the associated cyclic
\'etale cover have two-rank zero. All characters of
smaller conductor may survive.
\end{claim}
\begin{proof}
We first state and prove the field estimate used below.
If a subfield $M$ of a curve function field $F/k$ over
an algebraically closed field is generated by finitely
many functions each of degree $d$, and $e=[F:M]$,
then
\[
                    e\mid d,\qquad
                    g(M)\le(d/e-1)^2.                    \tag{AP.1}
\]
This is the controlled-field assertion of
\citet[Lemma 2.2 and Corollary 2.3, author p.~3]
{A_PrimPoonen2007}. For one generator it is immediate.
Suppose at an induction stage that $[F:M]=d/s$
and $g(M)\le(s-1)^2$. Adjoin another degree-$d$
function $h$, and put $a=[M(h):M]$. Then
\[
       [F:M(h)]=d/(as),\qquad [M(h):k(h)]=as.
\]
Castelnuovo--Severi gives
\[
 \begin{aligned}
 g(M(h))
 &\le a(s-1)^2+(a-1)(as-1)\\
 &=(as-1)^2-as(a-1)(s-1)\le(as-1)^2.
 \end{aligned}
\]
This proves the induction. In particular, the estimate
is valid for a compositum with many generators.

If $n$ is a prime power, A.C2.8 already proves the
claim. Suppose therefore that $r=\omega(n)\ge2$,
and put $d=2^r$. We have
\[
                            n>(d-1)^2.                   \tag{AP.2}
\]
For $r=2$, use $n\ge15>9$. For $r\ge3$,
$n\ge3\cdot5^{r-1}>4^r>(2^r-1)^2$:
the middle inequality holds at $r=3$ and its ratio
increases by a factor $5/4$ at each step.

Let $D\to C$ be the connected cyclic \'etale cover
associated with $L$, and let $G=\langle\sigma\rangle$
be its group. Its genus is $n+1$. The orbit hypothesis
and A.C2.5 kill all $\varphi(n)$ primitive characters,
so
\[
 f(D)\le f(C)+n-1-\varphi(n)
       \le n+1-\varphi(n)<n-1,                            \tag{AP.3}
\]
since $\varphi(n)\ge8$.
Define
\[
 a_\ell=\sigma^{n/\ell}\quad(\ell\mid n\text{ prime}),
 \qquad A=\prod_{\ell\mid n}(1-a_\ell).
                                                               \tag{AP.4}
\]
Its scalar on character $i$ is nonzero precisely when
$\gcd(i,n)=1$. The integral character argument of
A.C2.17 therefore makes $A$ zero on all geometric
two-primary torsion: over the finite unramified
extension of $\Z_2$ containing the $n$-th roots of
unity, primitive Tate summands have rank zero and
each other summand is killed by one factor.
Duality may invert indices but preserves primitivity.

A hypothetical collision of distinct $P,Q\in C(K)$
lifts after finite extension, by A.3, to $P',Q'\in D$
with two-primary difference. Applying $A$ gives the
principal divisor
\[
 \mathcal D=\sum_{S}(-1)^{|S|}
                    (a_SP'-a_SQ'),\qquad
 a_S=\prod_{\ell\in S}a_\ell,                             \tag{AP.5}
\]
where $S$ runs over subsets of the distinct primes
dividing $n$.
The $2^r$ elements $a_S$ are distinct, since their
prime-order factors have pairwise coprime orders.
The two deck orbits are disjoint and free.
Thus (AP.5) has $d$ distinct zeros and $d$ distinct
poles, all simple. Let $f$ be its rational function;
it is separable and has degree $d$.

Over $k=\overline K$ set
\[
              M=k(\gamma f:\gamma\in G),\qquad
              e=[k(D):M],
\]
and let $Z$ be the smooth curve with function field
$M$. Every generator has degree $d$, so (AP.1)
applies. The equality $e=1$ would give
$n+1=g(D)\le(d-1)^2$, contradicting (AP.2).
Hence $e>1$ and is a power of two.

The field $M$ is $G$-stable. The induced action on
$Z$ is faithful: its kernel has order dividing
both $e$ and the odd number $n$. It is free as
well. A nontrivial stabilizer contains an element
of odd prime order $\ell$; its free action on
$D$ partitions the degree-$e$ fibre divisor into
orbits of size $\ell$ with constant multiplicity.
Thus $\ell\mid e$, a contradiction.
It follows that
\[
                 g(Z)-1=n(g(Z/G)-1).
\]
The genus of $Z$ is therefore either one or at
least $n+1$. The bound (AP.1) is less than $n+1$,
by (AP.2), so $g(Z)=1$.

Write $f=h\circ v$ with $v:D\to Z$ of degree $e$.
Choose a zero of the nonconstant function $h$.
Its full inverse image is reduced by (AP.5), so
the zero is simple, the map $v$ is unramified
above it, and the fibre has exactly $e$ geometric
points. Each of $GP'$ and $GQ'$ maps injectively
to $Z$, because an equality of images would
give a stabilizer in its free $G$-action.
The fibre is contained in these two orbits, so
$e\le2$. Thus $e=2$.

The map $D\to Z$ is consequently a separable
quadratic cover of a genus-one curve. Its
involution $\tau$ commutes with $G$, because
$M$ is $G$-stable and the order-two deck group
has no nontrivial automorphism.
Its nonempty fixed set has cardinal divisible
by $n$. Riemann--Hurwitz gives total different
$2n$, and every wild quadratic fixed point
contributes at least two. There are exactly
$n$ fixed points. The ramified
Deuring--Shafarevich formula gives
\[
                         f(D)=2f(Z)+n-1\ge n-1,
\]
contradicting (AP.3)
\citep[Corollary 1.8, pp.~35--36]{A_PrimCrew1984}.
Alternatively, the independent classes
$n[R_i-R_0]$ from the ramification orbit, as
proved in A.C2.17, give the same lower bound
without the formula.

This excludes the collision. Finite-extension
descent and the open-subcurve assertion follow
from A.3 and functoriality to the compactification,
as in A.C2.17; no specialization of the theta
condition is required.

The new point is the support of (AP.5) in only
two deck orbits. Its degree is $2^{\omega(n)}$,
not necessarily four. The controlled field
and reduced zero fibre force the same quadratic
elliptic obstruction for every odd conductor.
This remains a conditional criterion: a single
theta torsion point need not meet all primitive
doubling orbits.
The certificate
\path{code/BC_arbitrary_primitive_prym_checks.py}
and its JSON output record 830 passing checks,
including 500 odd orders and six exact
cyclotomic identities at conductors with at
least three prime divisors.
\end{proof}

\subsection{An odd cyclic primitive factor in every genus and characteristic}

\begin{claim}{A.C2.20}{Proved}
Let $K$ be finitely generated over $\F_p$, and let
$C/K$ be smooth proper geometrically connected of
genus $g\ge2$. Suppose that geometrically there is
a connected cyclic \'etale cover $D\to C$ of odd
degree $n>1$ prime to $p$. For a deck generator
$\sigma$, put
\[
 P_n=\operatorname{im}(A)\subset J_D,\qquad
                 A=\prod_{\ell\mid n}(1-\sigma^{n/\ell}),
                                                               \tag{GP.1}
\]
where the product runs over distinct prime divisors.
If $f_p(P_n)=0$, then the geometrically prime-to-$p$
point-section map of $C$ is injective.
The conclusion holds after every finite extension,
for forms satisfying this geometric hypothesis,
and for hyperbolic open subcurves.

The image in (GP.1) is the primitive conductor-$n$
factor up to isogeny: its scalar on character $i$
is nonzero exactly when $\gcd(i,n)=1$.
Its isogeny class is independent of the generator.
In genus two and characteristic two, A.C2.5
identifies this hypothesis with the primitive
doubling-orbit condition in A.C2.19.
No theta characterization is asserted here in
higher genus.
\end{claim}
\begin{proof}
Suppose that distinct $P,Q\in C(K)$ give equal
point-section classes. After a finite extension,
define $D$, its action, and a chosen lift $P'$.
By A.3 there is a lift $Q'$ with
$[P'-Q']\in J_D[p^\infty]$.
The image of this class under $A$ belongs to
$P_n[p^\infty]$, which is zero by hypothesis.

Write $r=\omega(n)$, $d=2^r$, and
$a_S=\prod_{\ell\in S}\sigma^{n/\ell}$ for a
subset $S$ of the distinct primes dividing $n$.
Thus
\[
                 \sum_S(-1)^{|S|}(a_SP'-a_SQ')            \tag{GP.2}
\]
is principal. The prime-order factors have
pairwise coprime orders, so the $a_S$ are distinct.
The two deck orbits are free and disjoint.
Consequently (GP.2) has $d$ distinct zeros and
$d$ distinct poles, each simple. Its rational
function $f$ is separable of degree $d$.

For every odd $n>1$ one has
\[
                              n>(d-1)^2.                 \tag{GP.3}
\]
For $r=1$, use $n\ge3>1$; for $r=2$, use
$n\ge15>9$; for $r\ge3$, use
$n\ge3\cdot5^{r-1}>4^r>(2^r-1)^2$.
Let
\[
 M=\overline K(\gamma f:\gamma\in G),\qquad
                    e=[\overline K(D):M],
\]
and let $Z$ be its smooth proper curve.
The controlled-field induction proved in A.C2.19,
which uses only Castelnuovo--Severi, gives
\[
                       e\mid d,\qquad
                       g(Z)\le(d/e-1)^2.                 \tag{GP.4}
\]
It is valid in every characteristic
\citep[Proposition 2.1, Lemma 2.2 and Corollary 2.3,
author p.~3]{A_PrimPoonen2007}.
Since $g(D)=n(g-1)+1\ge n+1$, (GP.3) excludes
$e=1$. Thus $e>1$ is a power of two.

The field $M$ is $G$-stable. Its induced action
is faithful, since the kernel order divides both
$e$ and the odd number $n$. It is free: any
odd-prime-order stabilizer would partition the
degree-$e$ fibre divisor in $D$ into free orbits
with constant multiplicity, forcing that prime
to divide $e$.
Hence $g(Z)-1=n(g(Z/G)-1)$. Together with
(GP.3)--(GP.4), this forces $g(Z)=1$.

Factor $f=h\circ v$ through $v:D\to Z$.
The fibre above a zero of $h$ consists of
exactly $e$ unramified points, since the zero
divisor (GP.2) is reduced. Each of the two
original deck orbits maps injectively to $Z$,
by its free $G$-action. The whole zero fibre
lies in these two orbits, so $e\le2$. Thus
$e=2$. The resulting separable quadratic
involution $\tau$ commutes with $G$, since
$M$ is $G$-stable. For $r=1$, (GP.4) would
already contradict $g(Z)=1$; the following
argument covers every hypothetical quadratic
quotient obtained this way.

Put $E_0=Z/G$. The commuting actions give
$G\times\langle\tau\rangle$ on $D$, and
\[
 \overline K(E_0)
       =\overline K(D)^{G\times\langle\tau\rangle}.
\]
Thus $E_0$ has genus one and $C=D/G\to E_0$
is separable quadratic. It is ramified,
because $g(C)\ge2$. The map $Z\to E_0$ is
an odd \'etale isogeny after choosing origins.

For every $m\mid n$ put
$H_m=\langle\sigma^m\rangle$,
$D_m=D/H_m$, and $Z_m=Z/H_m$. The square
\[
 \begin{array}{ccc}
 D_m&\longrightarrow&Z_m\\
 \downarrow&&\downarrow\\
 C&\longrightarrow&E_0
 \end{array}                                              \tag{GP.5}
\]
has quadratic horizontal maps and degree-$m$
\'etale vertical maps. It is Cartesian.
Indeed $W=C\times_{E_0}Z_m$ has $W\to C$
\'etale, as the base change of $Z_m\to E_0$,
and hence $W$ is smooth. The extensions of
degrees two and $m$ over $\overline K(E_0)$
are linearly disjoint. Thus $W$ is integral,
and their compositum, of degree $2m$, is
$\overline K(D_m)$. This identifies $W$ with
$D_m$.

Suppose first that $p=2$. Let $r_1>0$ be
the number of geometric ramification points
of $C\to E_0$. By (GP.5), the corresponding
number for $D_m\to Z_m$ is $m r_1$:
every branch point has $m$ distinct preimages
under the \'etale isogeny and the same local
quadratic extension after base change.
All the genus-one curves $Z_m$ have the
same two-rank $\epsilon$, by isogeny.
The ramified Deuring--Shafarevich formula gives
\[
                       f_2(D_m)=m r_1+2\epsilon-1
                                                               \tag{GP.6}
\]
\citep[Corollary 1.8, pp.~35--36]{A_PrimCrew1984}.
Here $r_1$ is positive by Riemann--Hurwitz,
not an expression whose sign is left to
be determined.

Let $b_j$ be the rank of the exact
conductor-$j$ part of the geometric \'etale
two-adic Tate representation of $J_D$.
The odd-order averaging idempotent, together
with pullback and norm, identifies its
$H_m$-invariants with the Tate module of
$J_{D_m}$. Therefore
\[
                          f_2(D_m)=\sum_{j\mid m}b_j.
\]
The character scalars in (GP.1) identify
$b_n=f_2(P_n)$. M\"obius inversion yields
\[
 \begin{aligned}
 f_2(P_n)
 &=\sum_{m\mid n}\mu(n/m)f_2(D_m)\\
 &=r_1\sum_{m\mid n}\mu(n/m)m
   +(2\epsilon-1)\sum_{m\mid n}\mu(n/m)\\
 &=r_1\varphi(n)>0,
 \end{aligned}                                             \tag{GP.7}
\]
contradicting the hypothesis. No
one-dimensional character-space assumption
was used.

If $p$ is odd, the derived quadratic map
$C\to E_0$ is instead a geometrically
prime-to-$p$ Galois cover of a genus-one
curve. A.29l already gives point-section
injectivity for exactly this class, including
its finite-extension and finitely generated
form scopes. It excludes the assumed
collision. No Deuring--Shafarevich formula
for a group of order prime to $p$ is invoked.

The cover and quotient data descend after
finite extension, equality of section classes
restricts to that extension, and A.3 supplies
the required lifted $p$-primary difference.
The geometric contradiction therefore
establishes the assertion over $K$ and over
every finite extension. Forms are included
by the same reasoning, and equality on an
open curve maps to equality on its
compactification.

This is a sufficient geometric condition,
not an existence theorem for primitive
factors of zero $p$-rank. Nonordinariness
of a cover is weaker and does not imply it.
The certificate
\path{code/BC_general_primitive_prym_checks.py}
and its JSON output verify the M\"obius
identity symbolically at 250 odd orders and
in 56 numerical rank examples. The proof
above gives the general identity and its
geometric meaning.
\end{proof}


\subsection{A cyclic theta difference bound and forbidden conductors}

\begin{claim}{A.C2.21}{Proved}
Let $C$ be any smooth proper genus-two curve over
an algebraically closed field of characteristic two,
let $\Theta$ be its symmetric Raynaud theta curve,
and let $L$ have exact odd order $n>1$. Put
\[
 S=\{i\in\Z/n:L^i\in\Theta\},\qquad
 b(n)=\left\lfloor\frac{1+\sqrt{8n-7}}2\right\rfloor,
 \qquad \nu=\operatorname{ord}_n(2).
\]
Then
\[
                  |S|(|S|-1)\le2(n-1),\qquad |S|\le b(n).
                                                               \tag{C2.61}
\]
For the cyclic \'etale cover $D\to C$ associated with $L$,
let $f_{\rm prim}(D)$ denote the two-rank of its
primitive conductor-$n$ Jacobian factor. One has
\[
          f_{\rm prim}(D)\ge
                    \max\{0,\varphi(n)-\nu b(n)\}.             \tag{C2.62}
\]
If $-1\in\langle2\rangle\subset(\Z/n)^\times$, then
\[
          f_{\rm prim}(D)\ge
              \max\{0,\varphi(n)-\nu\lfloor b(n)/2\rfloor\}.
                                                               \tag{C2.63}
\]
In particular every cyclic \'etale cover of degree $127$
has primitive two-rank at least $14$, and every such
cover of degree $257$ has primitive two-rank at least
$80$. These two conductors cannot supply a primitive
factor of two-rank zero on any genus-two curve in
characteristic two. Their geometric deck actions on
two-primary torsion are faithful. These are method
obstructions and do not assert any point-section collision.
\end{claim}
\begin{proof}
For any nonzero degree-zero class $T$ there are
at most two ordered pairs $(P,Q)$ of geometric
points of $C$ with $[P-Q]=T$.
Indeed two distinct representations give
\[
                         P+Q'\sim P'+Q.
\]
The effective divisors cannot be formally equal
unless the pairs coincide or $P=Q$; the latter
would give $T=0$. Their distinct equivalence
therefore supplies the unique genus-two
hyperelliptic pencil. If $\iota$ is its
involution, the second pair must be
\[
                      (P',Q')=(\iota Q,\iota P).
\]
It is uniquely determined. This argument includes
ramification points and requires no transverse
intersection hypothesis.

The Abel embedding $R\mapsto\mathcal O(R-B_\theta)$
is injective. For each nonzero $k\in\Z/n$, an
ordered pair of distinct indices $i,j\in S$
with $i-j=k$ gives a distinct pair of curve
points whose difference is $L^k\ne0$.
There are at most two pairs for each $k$.
Summing over the $n-1$ nonzero differences proves
(C2.61), and solving the quadratic inequality
gives the stated exact integer value of $b(n)$.

Every primitive doubling orbit has length $\nu$.
By A.C2.5 its one-dimensional Frobenius cycle
contributes either its entire length or zero,
according as it avoids or meets theta.
Thus
\[
 f_{\rm prim}(D)=\nu\,
   \#\{\text{primitive doubling orbits missing }S\}.
                                                               \tag{C2.64}
\]
Each killed primitive orbit needs a distinct
index of $S$, so at most $b(n)$ can be killed.
There are $\varphi(n)/\nu$ primitive orbits,
which proves (C2.62).
Nonprimitive indices, including zero when it
belongs to theta, use part of the cardinality
bound without killing any primitive orbit.

If $-1$ is a power of two, negation preserves
each primitive doubling orbit. Theta symmetry
then makes every killed orbit contain both
$i$ and $-i$. These are distinct because $n$
is odd and $i$ is a unit. At most
$\lfloor b(n)/2\rfloor$ primitive orbits can
therefore be killed, proving (C2.63).

At $127$, $\varphi(127)=126$, $\nu=7$, and
$-1\notin\langle2\rangle$. The bound gives
$b(127)=16$, hence
\[
                    f_{\rm prim}(D)\ge126-7\cdot16=14.
\]
At $257$, one has $\nu=16$ and $2^8=-1$,
while $b(257)=23$. Consequently
\[
              f_{\rm prim}(D)\ge
                         256-16\lfloor23/2\rfloor=80.
\]
At prime conductor a surviving nontrivial
character is faithful. The integral odd-order
Tate decomposition in A.C2.8 therefore gives
faithful deck action on geometric two-torsion
and on the full two-adic Tate module.

The exact integer certificate
\path{code/A_cyclic_theta_difference_checks.py}
and its JSON output check $500$ odd orders,
using integer square roots and direct cycle
computations, and verify both displayed lower
bounds. The complete independent review is
\path{audit/A_cyclic_theta_difference_topos_review.md}.
The proof for all $n$ is the geometric difference
argument and (C2.64), not an extrapolation from
those finite checks. A.C2.14 separately rules
out one fixed finite prime support for all
curves. Neither obstruction excludes finding
a different suitable conductor for a particular
curve, and universal A1 remains A.30.
\end{proof}

\subsection{Forty-six ordinary genus-two moduli classes over \texorpdfstring{$\F_4$}{F4}}

\begin{claim}{A.C2.22}{Proved}
Among the $48$ geometric ordinary genus-two moduli
classes fixed by fourth-power Frobenius, the $46$
classes in the table below satisfy A1.
Equivalently, every ordinary genus-two curve over
$\F_4$ is covered by this conclusion unless its
normalized unordered coefficient triple has one
of the two symmetric polynomials
\[
 Z^3+\alpha Z^2+\alpha^2,\qquad
 Z^3+\alpha^2Z^2+\alpha,\qquad \alpha^2+\alpha+1=0.              \tag{C2.65}
\]
The assertion includes every finite extension,
every finitely generated form of the $46$
positive geometric classes, and every hyperbolic
open subcurve. The two classes in (C2.65) are not
decided by this enumeration; A.C2.25 below settles
both by a different complete primitive theta packet.
\end{claim}
\begin{proof}
\proofparagraph{The exact list of geometric classes.}
Over an algebraic closure every ordinary genus-two
curve in characteristic two has a reduced form
\[
 C(a,b,c):\quad y^2+y=\phi(x)=ax+\frac bx+\frac c{x+1},
                        \qquad abc\ne0.                        \tag{C2.66}
\]
Indeed the Artin--Schreier two-rank formula gives
three branch points, and the genus-two different
formula forces all three reduced pole orders
to be one. Normalize their coordinates to
$0,1,\infty$; a constant is removed after a finite
extension. These are the same normalizations
used in A.C2.1--A.C2.18.

Any geometric isomorphism preserves the unique
hyperelliptic pencil and permutes those three
branch coordinates. The changes $x\mapsto x+1$
and $x\mapsto1/x$ generate the permutations
of $a,b,c$, up to a constant.
Over the identity base map a generator change
has the form $y\mapsto e(x)y+v(x)$.
The coefficient of $y$ requires $e^2+e=0$,
so $e=1$. A nonconstant $v$ has a pole, and
$v^2+v$ then has even positive pole order,
which cannot equal a rational function with
only simple poles. Thus $v$ is constant.
Consequently the unordered triple $\{a,b,c\}$
is exactly the geometric moduli datum.

For a curve over $\F_4$, fourth-power Frobenius
preserves this triple. The possible triples
are precisely the root multisets of monic
cubics over $\F_4$ with nonzero constant term.
Their number is $4\cdot4\cdot3=48$.
There are ten split triples, eighteen triples
with splitting type $1+2$, and twenty irreducible
cubic triples. Their coefficients lie respectively
in $\F_4,\F_{16},\F_{64}$.
Repeated coefficient roots are allowed; the
three branch coordinates remain distinct.
Only containment of every actual $\F_4$ curve
in this list is needed, so no separate descent
existence theorem for stable moduli points is
being assumed.

Encode $\F_4$ by $0,1,2=\alpha,3=\alpha^2$.
A key $s_1s_2s_3$ means $Z^3+s_1Z^2+s_2Z+s_3$.
In the following table $q$ is the coefficient-field
size, $n$ is a theta-class annihilator, and $e$
is the degree of the chosen coordinate factor
over that field. The omitted keys are exactly
$203,302$.
\begin{center}
\begin{tabular}{cccc@{\qquad}cccc}
\toprule
Key & $q$ & $n$ & $e$ & Key & $q$ & $n$ & $e$\\
\midrule
001 & 4 & 7 & 1 & 133 & 4 & 15 & 2\\
002 & 64 & 27 & 1 & 201 & 64 & 3 & 1\\
003 & 64 & 27 & 1 & 202 & 16 & 23 & 1\\
011 & 64 & 5 & 1 & 211 & 16 & 3 & 1\\
012 & 16 & 7 & 1 & 212 & 4 & 9 & 2\\
013 & 16 & 7 & 1 & 213 & 64 & 5 & 1\\
021 & 64 & 13 & 1 & 221 & 16 & 3 & 1\\
022 & 16 & 15 & 1 & 222 & 64 & 9 & 1\\
023 & 16 & 5 & 1 & 223 & 4 & 7 & 2\\
031 & 64 & 13 & 1 & 231 & 4 & 3 & 1\\
032 & 16 & 5 & 1 & 232 & 64 & 19 & 1\\
033 & 16 & 15 & 1 & 233 & 64 & 11 & 1\\
101 & 64 & 3 & 1 & 301 & 64 & 3 & 1\\
102 & 16 & 49 & 3 & 303 & 16 & 23 & 1\\
103 & 16 & 49 & 3 & 311 & 16 & 3 & 1\\
111 & 4 & 3 & 1 & 312 & 64 & 5 & 1\\
112 & 64 & 7 & 1 & 313 & 4 & 9 & 2\\
113 & 64 & 7 & 1 & 321 & 4 & 3 & 1\\
121 & 16 & 3 & 1 & 322 & 64 & 11 & 1\\
122 & 4 & 15 & 2 & 323 & 64 & 19 & 1\\
123 & 64 & 11 & 1 & 331 & 16 & 3 & 1\\
131 & 16 & 3 & 1 & 332 & 4 & 7 & 2\\
132 & 64 & 11 & 1 & 333 & 64 & 9 & 1\\
\bottomrule
\end{tabular}
\end{center}
The distinct displayed annihilators are
\[
                  3,5,7,9,11,13,15,19,23,27,49.
\]
For every one of them, direct multiplication
shows $\langle2,-1\rangle=(\Z/n)^\times$.
This finite table is a list of positive
certificates, not a conclusion from negative
searches at other orders.

\proofparagraph{The fixed exact witnesses.}
The complete data in
\path{code/A_theta_F4_moduli_witnesses.json}
give, for every row, the exact coefficient
triple, an irreducible coordinate polynomial
$q_0(T)$ over its coefficient field, a root
$t\ne0,1$, and $\delta=\phi(t)$.
They also give polynomials $A(x),B(x)$ with
\[
 \deg A\le d,\quad B\text{ monic of degree }d=(n-3)/2,
 \quad A(t)=0,\quad B(t)B(0)B(1)\ne0.
\]
With $u=x(x+1)$ and
$F=b+(b+c)x+ax^2+ax^3=u\phi(x)$, each record
verifies the full polynomial identity
\[
            u(A^2+AB)+B^2(F+\delta u)=a(x+t)^n.                \tag{C2.67}
\]
The coefficient fields have respective binary
moduli $7,19,67$, namely
\[
 w^2+w+1,\qquad w^4+w+1,\qquad w^6+w+1.
\]
Elements of the coordinate field are encoded
by consecutive blocks of $2,4,$ or $6$ bits
for their coefficients in $1,T,\ldots,T^{e-1}$.
All chosen coordinate factors have degree
at most three. The largest coordinate fields
are $\F_{4096}$ in the two degree-$49$ rows.
The auxiliary $\beta$ below may require a
further quadratic constant extension.

The exact verifier
\path{code/A_theta_F4_moduli_witnesses.py}
uses ordinary polynomial coefficient lists
and imports none of the determinant engine.
It checks the finite-field operations and
irreducibility, solves the original local
jet equations, verifies all those jets, and
multiplies (C2.67) in full. Every one of the
$46$ fixed records passes. Thus the positive
certificates can be checked without any
large determinant or inference from numerical
point counts.

Choose $\beta^2+\beta=\delta$ and put
$R=(t,\beta)$, $h=A+B(y+\beta)$.
Its norm is $a(x+t)^n/u$. At $R$, $h=0$
and its conjugate has nonzero value $B(t)$.
The $x$-map is \'etale there, and the norm
has its only zero at $t$, of order $n$.
The values $B(0),B(1)\ne0$ give simple poles
at $W_0,W_1$. At infinity, $x^dy$ has the
unique largest pole order $2d+1=n-2$.
Therefore
\[
 \begin{aligned}
 \operatorname{div}(h)&=nR-W_0-W_1-(n-2)W_\infty,\\
 \operatorname{div}(h/u^{(n-1)/2})&=n(R-B_\theta).
 \end{aligned}                                                 \tag{C2.68}
\]
The theta class $L=\mathcal O(R-B_\theta)$
is nonzero: otherwise $2R$ is canonical,
forcing $R$ to be Weierstrass, contrary to
$t\ne0,1,\infty$.
Its actual order $n_0>1$ divides $n$.
The reduction map $(\Z/n)^\times\to(\Z/n_0)^\times$
is surjective, so the generation property
of $2,-1$ descends to $n_0$.
The compact criterion in A.C2.17, or A.C2.19,
proves A1. In particular, exactness of the
listed annihilator is unnecessary, and no
negative theta scan enters this implication.
The cover, isomorphism and parameters descend
after finite extension of the field of
definition; A.3 and the geometric contradiction
then give all stated form and extension
consequences. Equality on an open subcurve
maps to equality on its compactification.

\proofparagraph{The largest two witnesses are small finite-field identities.}
For key $102$, take $\F_{16}$ with
$w^4+w+1=0$, $w$ encoded by $2$, and
$(a,b,c)=(6,8,15)$. Its coordinate polynomial is
\[
                         q_0(T)=T^3+2T^2+4T+14.
\]
For key $103$, take $(a,b,c)=(7,10,12)$ and
\[
                         q_0(T)=T^3+5T^2+2T+13.
\]
Neither cubic has a root in $\F_{16}$, so both
are irreducible. The following matrices list
the coefficients of $A$ and $B$ in increasing
powers of $x$, read row by row; integers use
four-bit coefficient blocks in the indicated
$\F_{4096}$.
For key $102$, $\delta=1663$, and
\[
\begin{aligned}
\operatorname{coeff}(A)&=\left[\begin{matrix}
2967 & 3094 & 2226 & 1784 & 3122 & 1282\\
2330 & 2961 & 160 & 1454 & 3913 & 1774\\
2664 & 2041 & 396 & 120 & 2115 & 3265\\
3418 & 2933 & 2951 & 2846 & 2919 & 1547
\end{matrix}\right],\\[3pt]
\operatorname{coeff}(B)&=\left[\begin{matrix}
1357 & 1974 & 2300 & 2749 & 1451 & 1952\\
3656 & 3667 & 1719 & 3593 & 4043 & 3650\\
2204 & 3665 & 3643 & 1821 & 57 & 3133\\
1148 & 362 & 3621 & 1909 & 111 & 1
\end{matrix}\right].
\end{aligned}
\]
For key $103$, $\delta=1900$, and
\[
\begin{aligned}
\operatorname{coeff}(A)&=\left[\begin{matrix}
3766 & 2071 & 2789 & 1994 & 2117 & 773\\
2847 & 3761 & 240 & 1021 & 3115 & 2013\\
3962 & 1739 & 424 & 106 & 2596 & 2177\\
2367 & 3683 & 3750 & 3613 & 3702 & 1806
\end{matrix}\right],\\[3pt]
\operatorname{coeff}(B)&=\left[\begin{matrix}
809 & 1767 & 2760 & 4073 & 1022 & 1776\\
3370 & 3380 & 2022 & 3339 & 3214 & 3365\\
2744 & 3377 & 3406 & 1561 & 75 & 2121\\
616 & 383 & 3411 & 1635 & 124 & 1
\end{matrix}\right].
\end{aligned}
\]
These two identities alone verify the largest
positive entries without recomputing their
$25$ maximal minors. The separate discovery
scan also excluded theta seven-torsion there,
so the displayed class actually has exact
order $49$; this stronger fact is not needed
for the table's point-section conclusion.

\proofparagraph{The actual remaining pair.}
For key $203$, the normalized coefficients are
\[
 \F_{16}=\F_2[w]/(w^4+w+1),\quad \alpha=w^2+w,
 \qquad(a,b,c)=(1,w^3+1,w^3+w^2+w)=(1,9,14).
\]
Fourth-power Frobenius swaps $b,c$, inducing
branch permutation $0\leftrightarrow1$.
An explicit $\F_4$ model is
\[
              v^2+v=z+\frac{\alpha^2(z+1)}{z^2+z+\alpha}.
                                                               \tag{C2.69}
\]
Indeed put $x=z+w$ and $y=v+k$ with
$k^2+k=w$, which has a solution in $\F_{16}$.
The rational function transforms exactly
as displayed. Key $302$ is the coefficientwise
square conjugate, with $\alpha$ replaced by
$\alpha^2$ and normalized coefficients
$(1,11,13)$.
These two classes have not been decided by
the table. Bounded further scans and partial
primitive-orbit hits are kept separate from
this positive theorem; they yield neither
an A1 counterexample nor a universal bound.

The exact geometric discovery criterion is
proved under A.C2.18 and in
\path{audit/A_theta_Hankel_scan.md}.
Original separate-determinant records are
preserved in
\path{code/A_theta_F4_moduli_original_minors.json}.
An augmented-Bareiss adjugate implementation
was compared with $189$ identical exact gcds
from those records; the check is
\path{code/A_theta_F4_moduli_adjugate_check.json}.
The fixed norm certificates are an independent
positive verification and are sufficient for
the present theorem.
A wholly independent SymPy quotient-ring implementation,
\path{code/A_F4_norm_root_independent.py}, also checks
all $46$ identities, field and coordinate irreducibility,
local values, moduli keys and unit-group generation;
its adjacent JSON records PASS. The complete proof and
scope record are
\path{audit/A_F4_moduli_46_candidate.md}.
Universal A1 remains A.30.
\end{proof}


\subsection{Descent of a primitive-zero factor through an abelian total cover}

\begin{claim}{A.C2.23}{Proved}
Let $k$ be algebraically closed of characteristic $p>0$,
let $C/k$ be a smooth proper connected curve, and let
$W\to C$ be a connected finite \'{e}tale Galois cover
with finite abelian group $A$, of odd order prime to $p$.
Let $H\subset A$ be a nontrivial cyclic subgroup of
order $n$, let $h$ generate $H$, and put $V=W/H$.
Suppose that
\[
 P_H=\operatorname{im}\left(
       \prod_{\ell\mid n}(1-h^{n/\ell}):J_W\longrightarrow J_W
                         \right)
\]
has $p$-rank zero. Then an intermediate quotient
$D=W/K$ is a cyclic \'{e}tale cover of $C$, of odd
order $N>1$ prime to $p$, whose primitive conductor-$N$
factor has $p$-rank zero. One can choose $N$ divisible
by $n$.

Thus the existence of a cyclic primitive-zero factor
over an intermediate base inside an abelian total cover
is equivalent to the existence of a cyclic
primitive-zero cover of the original curve.
For genus at least two over a finitely generated field,
this condition implies A1 by A.C2.20.
The assertion concerns an abelian \emph{total} cover
$W/C$; a sequence of abelian covers can have a
nonabelian normal closure over $C$ and is not
covered by this equivalence.
\end{claim}
\begin{proof}
\proofparagraph{Primitive characters and geometric Tate modules.}
Write
\[
 \mathcal T_W=T_p(J_W)\otimes_{\mathbf Z_p}
                              \overline{\mathbf Q}_p.
\]
The geometric $p$-adic Tate module has rank equal
to the $p$-rank. Since $A$ has order prime to $p$,
its action is semisimple; since it is abelian,
\[
 \mathcal T_W=\bigoplus_{\psi\in A^\vee}\mathcal T_{W,\psi},
 \qquad
 A^\vee=\operatorname{Hom}(A,\overline{\mathbf Q}_p^\times).
\]
On the $\psi$-space the operator defining $P_H$
has scalar
\[
              \prod_{\ell\mid n}
                   \bigl(1-\psi(h)^{n/\ell}\bigr).
\]
This scalar is nonzero precisely when $\psi|_H$
is faithful. Indeed, if its order $m$ properly
divides $n$, some prime $\ell$ satisfies
$v_\ell(m)<v_\ell(n)$, whence $m\mid n/\ell$.
If its order is $n$, no factor vanishes.

The image on a rational geometric Tate module of
a homomorphism of abelian varieties is the Tate
module of its image abelian subvariety. This
follows by factoring through the connected
kernel and the resulting isogeny; isogenies
become isomorphisms on rational $p$-adic Tate
modules. Consequently
\[
 \mathcal T_{W,\psi}=0
          \quad\text{if }\psi|_H\text{ is faithful}.
                                                        \tag{AT.1}
\]
Here zero $p$-rank is used, not merely a
nonzero Frobenius kernel.

\proofparagraph{Extension of a faithful subgroup character.}
Choose a faithful character
$\theta:H\to\overline{\mathbf Q}_p^\times$.
It extends to a character $\chi$ of $A$.
Here is the elementary extension argument.
Suppose a character is defined on a subgroup $B$,
choose $a\in A$, and let $r$ be the order of
$a+B$ in $A/B$. Choose $z$ with
$z^r=\theta(ra)$. The formula
\[
               b+ja\longmapsto\theta(b)z^j
\]
defines a character of $B+\langle a\rangle$.
An equality of two such expressions has
$j-j'$ divisible by $r$, and the root relation
gives precisely the necessary equality of
values. If $a$ has order $s$, then $r\mid s$
and $z^s=\theta(sa)=1$, so the image remains
finite. Finitely many enlargements give $\chi$.

Let $N$ be the order of $\chi$, put
$K=\ker\chi$, and set $D=W/K$.
Then $A/K$ is cyclic of order $N$ and
$D\to C$ is a connected cyclic \'{e}tale cover.
Since $\chi|_H$ has order $n$, one has $n\mid N$.
For every $u\in(\mathbf Z/N)^\times$,
\[
                    (\chi^u)|_H=\theta^u
\]
is faithful: $n\mid N$ implies $\gcd(u,n)=1$.
Equation (AT.1) therefore gives
\[
       \mathcal T_{W,\chi^u}=0
                 \quad(u\in(\mathbf Z/N)^\times).
                                                        \tag{AT.2}
\]

\proofparagraph{Norm and pullback for the cyclic quotient.}
For $q:W\to D$, pullback and norm satisfy
\[
 q_*q^*=[|K|],\qquad
 q^*q_*=\sum_{\gamma\in K}\gamma^*.
\]
Since $|K|$ is prime to $p$, they identify
$T_p(J_D)$ with $T_p(J_W)^K$, already
integrally. For an invariant vector $v$, the
inverse comparison is explicit:
\[
                   v=q^*\bigl(|K|^{-1}q_*v\bigr).
\]
The $A$-characters trivial on $K$ are exactly
the powers of $\chi$. The faithful characters
of $A/K$ are exactly the primitive powers
$\chi^u$, $u\in(\mathbf Z/N)^\times$.
Their Tate spaces vanish by (AT.2), so the
primitive conductor-$N$ factor of $J_D$ has
$p$-rank zero.
Dualizing conventions may invert character
indices; faithfulness and primitivity are
unchanged.

This proves descent. The converse existence
implication takes $W=D$, $A=H$, and $V=C$.
All finite geometric covering data descend
after a finite extension of a field of
definition. Thus A.C2.20 supplies the stated
finitely generated-field consequence, including
its forms and open-subcurve variants.
\end{proof}

For genus two in characteristic two, this
explains an exact obstruction to an abelian
base-change strategy. Passing to an
intermediate base groups the original
character spaces according to their
restrictions to the new deck subgroup.
One theta hit need not kill all original
characters extending a faithful character
upstairs. If every such contribution does
vanish, the proof above already produces a
cyclic quotient over the original curve
with all primitive doubling orbits killed.
In particular A.C2.14's finite-prime-set
avoidance cannot be circumvented by taking
intermediate cyclic stages of an abelian
total cover supported on the same prime set.

This is a limitation of that construction,
not a counterexample to A1 or a
nonexistence assertion for nonabelian
normal closures. The separate source audit
\path{audit/A_nonabelian_Raynaud_gap.md}
records why the opened higher-rank
nonordinary-cover theorem does not itself
supply zero stable multiplicities or a
sparse integral annihilator.
The finite checks in
\path{code/BC_abelian_descent_and_support_checks.py}
verify 31,006 character assertions and
an explicitly nongeometric Heisenberg
module illustrating the distinction
between representation support and
multiplicity.


\subsection{Exact structural limits for the last two ordinary classes}

\begin{claim}{A.C2.24}{Proved}
Let $\alpha^2+\alpha+1=0$ and let
\[
 C_0/\mathbf F_4:\qquad
 v^2+v=z+\frac{\alpha^2(z+1)}{z^2+z+\alpha}.
                                                        \tag{C2.70}
\]
The following assertions hold for $C_0$ and its
coefficientwise-square conjugate, the two classes
left undecided by A.C2.22 and subsequently
settled by A.C2.25.
\begin{enumerate}
\item Their ordinary Jacobians are geometrically simple, with
\[
 \operatorname{End}^0_{\overline{\mathbf F}_4}(J_{C_0})
 =\mathbf Q(\pi)
 =\mathbf Q\bigl(\sqrt2,\sqrt{-7+4\sqrt2}\bigr),
                                                        \tag{C2.71}
\]
a non-Galois quartic CM field. Their geometric
automorphism group is exactly $C_2$. In particular
the original curves have no nonconstant elliptic
map or map to an abelian variety of two-rank zero.
\item Every geometric isogeny between $J_{C_0}$ and the
Jacobian of the quadratic Artin--Schreier twist of
the A.C2.13 split-branch curve has even degree.
\item The geometric Raynaud theta divisor has no
nonzero $\ell$-primary point for
\[
 \ell\in\{3,5,7,11,13,19,23,29,31,37,43,53,59,61,67\}.
                                                        \tag{C2.72}
\]
It has no point of order supported on any of
\[
 \{3,5\},\quad\{3,7\},\quad\{5,7\},\quad
 \{3,11\},\quad\{7,11\}.                              \tag{C2.73}
\]
Arbitrarily high powers of the listed primes are
included in these exclusions.
\item Its full geometric theta-$73$ locus consists of
four points on one cyclic subgroup. Relative to a
generator on theta, the packet is
\[
                    \{1,27,46,72\}=\{\pm1,\pm27\}.
                                                        \tag{C2.74}
\]
Only four of the eight primitive doubling orbits
meet theta. Every cyclic $73$-power cover therefore has
positive primitive two-rank; the degree-$73$ cover belonging
to this packet has primitive two-rank $36$ and
total two-rank $38$.
\end{enumerate}
These are exact limitations of specified
constructions. A1 for these two curves is not
decided by this claim.
\end{claim}
\begin{proof}
\proofparagraph{Counts and geometric simplicity.}
Direct Artin--Schreier enumeration on the smooth
normalization gives
\[
 \#C_0(\mathbf F_4)=7,\qquad
 \#C_0(\mathbf F_{16})=15,\qquad
 \#C_0(\mathbf F_{256})=255.
\]
At a simple rational pole there is one point;
elsewhere there are two precisely when the
absolute trace of the right side is zero.
The first two counts give
\[
           P(X)=X^4+2X^3+X^2+8X+16                     \tag{C2.75}
\]
for fourth-power Frobenius. The third count checks
its fourth-power trace. We use the exact
curve-to-Jacobian trace formula in
\cite[III, Theorem~11.1 and Corollary~11.4]{A_Milne2008},
together with the abelian-variety count and Weil
absolute values in its II, Theorem~1.1.
The three simple Artin--Schreier poles give
ordinariness.

Modulo three, $P$ is $\Phi_{10}$, irreducible
because $3$ has order four modulo $10$.
The squarefree nontrivial-root-ratio polynomial,
computed from $\operatorname{Res}_X(P(X),P(ZX))$
after deleting $(Z-1)^4$, is
\[
 \begin{split}
 R(Z)={}&64Z^8+80Z^7+108Z^6+7Z^5+26Z^4\\
        &\hspace{12mm}+7Z^3+108Z^2+80Z+64.
 \end{split}                                           \tag{C2.76}
\]
Its degree is at most eight for a structural
reason: write the four Weil roots as
$a,4/a,b,4/b$. There are four ratios within
reciprocal pairs and four cross ratios, the
latter each repeated twice. A root of unity
among these ratios would give a cyclotomic
factor of degree at most eight. The full list
of orders to check is
\[
 1,2,3,4,5,6,7,8,9,10,12,14,15,16,18,20,24,30.
\]
This list is exhaustive: only primes $2,3,5,7$
can divide such an order, with exponents at
most $4,2,1,1$, and the remaining finite
products give exactly the list. Exact
polynomial gcds with all eighteen cyclotomic
polynomials equal one.

Thus distinct roots of $P$ never become equal
after taking a positive power. Every $\pi^m$
has minimal polynomial of degree four, equal
to its irreducible characteristic polynomial.
A proper geometric abelian subvariety would
descend to some finite extension and factor
that polynomial. This proves geometric
simplicity.

\proofparagraph{The entire rational endomorphism algebra.}
Every geometric endomorphism commutes with
$\pi^m$ for some $m$. The tensor injection
\[
 \operatorname{End}(J)\otimes\mathbf Z_\ell
          \hookrightarrow\operatorname{End}(T_\ell J)
 \qquad(\ell\ne2)
\]
is \cite[I, Theorem~10.15, printed p.~49]{A_Milne2008}.
The centralizer of $\pi^m$ on the
four-dimensional rational Tate module is
$\mathbf Q_\ell[\pi^m]=\mathbf Q_\ell[\pi]$:
the operator has four distinct eigenvalues
and minimal polynomial of degree four.
All geometric endomorphisms therefore lie in
the same four-dimensional centralizer.
Their rational algebra contains the
degree-four field $\mathbf Q(\pi)$, so
equality follows. No assertion about a
maximal integral order is needed.

Put $s=\pi+4/\pi$. Equation (C2.75) gives
$s^2+2s-7=0$, and solving
$\pi^2-s\pi+4=0$ gives (C2.71).
Both real conjugates of $-7+4\sqrt2$ are
negative. Their product is $17$. If the CM
field were normal, it would contain both
square roots and hence $\sqrt{17}$.
Together with $\sqrt2$, this would generate
a totally real degree-four subfield, forcing
the whole field to be real, a contradiction.

An elliptic map would give an elliptic
Jacobian quotient. A nonconstant map to a
two-rank-zero abelian variety would give a
nonzero abelian image of the simple ordinary
Jacobian, hence an ordinary image of dimension
two. Both are impossible. Also the normalized
coefficients $(1,9,14)$ over
$\mathbf F_{16}=\mathbf F_2[w]/(w^4+w+1)$
are pairwise distinct. The six branch
permutations act by their coefficient
permutations, and an Artin--Schreier correction
cannot alter a simple-pole coefficient.
A.C2.22's classification therefore leaves
only the hyperelliptic involution.

\proofparagraph{Why a specific odd-isogeny transfer fails.}
The quadratic twist of the A.C2.13 curve
has the same polynomial (C2.75), as checked
independently by its counts $7,15,255$.
All three of its ramification points are
rational over $\mathbf F_4$, so Frobenius
fixes all three nonzero geometric two-torsion
classes. For $C_0$, Frobenius exchanges the
two finite ramification points and fixes
infinity. Its action on the three nonzero
classes is a transposition. Since $P(1)=28$,
the two rational two-primary groups are
respectively $C_2\times C_2$ and $C_4$.

For every $m$, the equality
$\mathbf Q(\pi^m)=\mathbf Q(\pi)$ gives
$R_m\in\mathbf Q[X]$ such that
$R_m(X^m)=X$ modulo $P(X)$.
The same polynomial works for both Jacobians.
A geometric homomorphism defined over
$\mathbf F_{4^m}$ intertwines their $m$-th
Frobenius powers; applying $R_m$ shows that
it intertwines their original fourth-power
Frobenius as well. An odd-degree isogeny
would give a Frobenius-equivariant isomorphism
of geometric two-torsion, contradicting the
two different actions. This concerns the
specified comparison pair, not all possible
isogeny correspondences or further covers.

\proofparagraph{Uniform torsion exponents from exact operator identities.}
If $P(X)\mid X^m-1$ modulo an odd prime
$\ell$, then $\pi^m$ fixes $J[\ell]$.
A.C2.12 makes every $\ell$-primary theta
point fixed by $\pi^m$. Suppose
\[
 (X^m-1)B(X)\equiv\ell^e
                  \pmod{P(X),\ell^{e+1}}.              \tag{C2.77}
\]
Because $P$ is monic, this yields an integral
endomorphism identity
\[
 (\pi^m-1)B(\pi)=\ell^e(1+\ell R_0(\pi)).
\]
On geometric $\ell$-primary torsion the
last factor is invertible, by its geometric
series at each finite torsion level. A
$\pi^m$-fixed point is therefore killed by
$\ell^e$. This does not presume that an
integral Frobenius lattice is cyclic.

Here are the exact certificates. The final
column lists the coefficients of $B$ in
ascending order; each row is checked directly
in the finite quotient ring in (C2.77).
\begin{center}
\scriptsize
\begin{tabular}{r r r l}
$\ell$&$m$&$\ell^e$&$B$ coefficients\\\hline
3&10&3&$(4,1,6,0)$\\
5&52&5&$(18,9,6,19)$\\
7&12&7&$(17,46,7,48)$\\
11&610&11&$(76,7,93,72)$\\
13&42&13&$(81,150,123,166)$\\
17&1224&289&$(1121,4202,390,3232)$\\
19&180&19&$(341,50,82,278)$\\
23&132&23&$(90,363,491,378)$\\
29&11788&29&$(425,801,273,143)$\\
31&120&31&$(379,567,530,11)$\\
37&24660&37&$(1116,23,450,215)$\\
41&140&1681&$(66658,64498,20991,27342)$\\
43&924&43&$(1238,65,550,591)$\\
47&23&2209&$(68200,85019,58876,86032)$\\
53&702&53&$(2277,11,505,1824)$\\
59&1740&59&$(851,1200,2335,1832)$\\
61&111660&61&$(3236,2301,2389,134)$\\
67&2244&67&$(201,915,1051,1020)$\\
71&1260&5041&$(157405,262215,277523,203927)$\\
73&1332&73&$(4521,1064,2697,1508)$
\end{tabular}
\end{center}
For every prime in (C2.72), the bound is
$\ell$ itself and the complete A.C2.18
Hankel test has residual maximal-minor gcd
one. This proves absence at every power.
The bounds at $17^2,41^2,47^2,71^2$ are
not asserted to be attained or excluded.

For a finite set $S$, first apply A.C2.14
to a common Frobenius power fixing every
$J[\ell]$, $\ell\in S$. It fixes the whole
$S$-supported theta point, hence every
primary component. Apply (C2.77) to these
components at that same power. The exact
component certificates give
\[
\begin{array}{c|r|r}
S&m&\text{order divides }N\\\hline
\{3,5\}&260&75\\
\{3,7\}&60&63\\
\{5,7\}&156&35\\
\{3,11\}&610&33\\
\{7,11\}&3660&77
\end{array}
\]
The complete geometric theta-$N$ test is
empty in each row. Thus (C2.73) includes
all powers, not only squarefree conductors.
The origin is outside the ordinary Raynaud
theta divisor. The exact ten component
polynomials and their remainders are in
\path{code/A_F4_remaining_mixed_packets.json}.

\proofparagraph{The complete theta-seventy-three packet.}
Use the binary encoding in A.C2.22, with
$w=2$ and $\alpha=w^2+w=6$ in $\mathbf F_{16}$.
The complete Hankel gcd is
\[
             T^2+6T+9=(T+8)(T+14).
\]
At these two coordinates the
Artin--Schreier constants are $9$ and $15$,
both of absolute trace one over $\mathbf F_{16}$.
Thus the four points lie over $\mathbf F_{256}$,
with sixteenth-power Frobenius exchanging
the two points above each coordinate.
The full norm witnesses verify order $73$;
alternatively A.C2.18's complete kernel
criterion gives annihilation by $73$, and
the non-Weierstrass condition makes the
class nonzero.

Equation (C2.75) gives
\[
                 \#J(\mathbf F_{256})=2^7\cdot7\cdot73.
\]
Its $73$-primary group is therefore cyclic
of order $73$, and contains all four theta
points. On normalized coordinates, actual
fourth-power Frobenius is
\[
                 (x,y)\longmapsto(x^4+1,y^4+\alpha).
\]
Its square is the hyperelliptic involution
on the four points. The theta origin class
$B_\theta$ is rational and satisfies
$2B_\theta=K_C$, so that involution acts by
negation on theta classes. Hence, if
$\pi\eta=k\eta$ for a theta generator,
$k^2=-1$ modulo $73$, giving $k=\pm27$.
The two coordinates are exchanged by
$\pi$, so the four points are exactly
$\pm\eta,\pm27\eta$.

The order of $2$ modulo $73$ is $9$.
The four multipliers in (C2.74) occupy
four distinct doubling orbits; four of
the eight remain. A.C2.5 gives primitive
two-rank $4\cdot9=36$, plus the base rank
two. Any other cyclic $73$ subgroup has
no theta point, so its primitive two-rank
is $72$. The exponent bound in the
$\ell=73$ row puts every $73$-primary
theta point in $J[73]$. Thus a cyclic
$73^a$-cover for $a\ge2$ has no primitive
theta character at all and has primitive
two-rank $\varphi(73^a)>0$.
This is a complete
$73$-primary statement, not merely a
rational-point sample. The separate Cantor
and scalar-Frobenius calculation reproduces
the same packet exactly.

\proofparagraph{Certificates and precise remaining scope.}
The author structural certificate is
\path{code/A_F4_remaining_structure.py};
the independent implementation
\path{code/A_F4_remaining_structure_degeneration.py}
uses its own finite-field and quotient-polynomial
arithmetic and checks the three counts,
all eighteen cyclotomic exclusions, and
all thirty primary and mixed identities.
The complete review is
\path{audit/A_F4_remaining_structure_degeneration_review.md}.
The complete conductor-$73$ gcd and norm
records are in
\path{code/A_theta_F4_remaining_73_89.json}
and
\path{code/A_F4_remaining_primitive_packets.py},
with the latter's adjacent JSON.
Root independently read the proof, reopened
the Milne passages, and reran the structural
and complete-packet certificates.

None of these assertions supplies colliding
sections. A.C2.25 below uses a different
four-prime conductor to settle A1 for both
curves, while every obstruction proved here
retains its stated scope.
\end{proof}


\subsection{A complete large-conductor packet and all ordinary genus-two classes over \texorpdfstring{$\F_4$}{F4}}

\begin{claim}{A.C2.25}{Proved}
Let $\alpha^2+\alpha+1=0$. The smooth proper curve
\[
 C_0/\F_4:\qquad
 v^2+v=z+\frac{\alpha^2(z+1)}{z^2+z+\alpha}                 \tag{C2.78}
\]
and its coefficientwise-square conjugate satisfy A1.
Consequently all $48$ geometric ordinary genus-two
moduli classes fixed by fourth-power Frobenius
satisfy A1. In particular this holds for every
ordinary genus-two curve over $\F_4$, over every
finite extension, for every finitely generated
form of any of these geometric classes, and
for every hyperbolic open subcurve.

The additional two classes are certified by a
connected cyclic \'etale cover of odd degree
\[
 M=7\cdot31\cdot2503\cdot4519=2\,454\,499\,369             \tag{C2.79}
\]
whose primitive conductor-$M$ Jacobian factor
has two-rank zero. Universal A1 remains A.30.
\end{claim}
\begin{proof}
\proofparagraph{The curve and the exact theta origin.}
Set
\[
 h=z^2+z+\alpha,\qquad G=z^3+z^2+z+\alpha^2,
 \qquad Y=hv,
\]
so that the polynomial model is
\[
                         Y^2+hY=f,\qquad f=hG.             \tag{C2.80}
\]
Here $f$ is monic of degree five. At a root
of $h$, the value $G=\alpha^2(z+1)$ is nonzero;
thus the affine model is smooth. The original
Artin--Schreier equation has three simple poles,
the two roots of $h$ and infinity, and defines
an ordinary genus-two curve. Write $W_1,W_2$
for its finite ramification points and $W_\infty$
for its point at infinity. The different and
the divisor of $dz$ give
\[
 \operatorname{div}(dz)=2(W_1+W_2-W_\infty),\qquad
 B_\Theta=\mathcal O(W_1+W_2-W_\infty).                     \tag{C2.81}
\]
The locally exact differential $dz$ identifies
this with the Raynaud bundle in A.C2.5. Hence
\[
 \Theta=\{[R-B_\Theta]:R\in C_0(\overline{\F}_2)\}.
\]
The pair $T=(h,0)$ represents
$[W_1+W_2-2W_\infty]$. Its double is zero,
since $\operatorname{div}(h)=2W_1+2W_2-4W_\infty$.
Consequently
\[
                  [R-W_\infty]+T=[R-B_\Theta].             \tag{C2.82}
\]
This fixes both the origin and the sign of the
theta test. The individual finite branch points
need not be rational over the certificate field;
their sum and $T$ are rational.

Choose $r$ with $r^2+r=\alpha$ and put $x=z+r$.
After removing a constant, (C2.78) has the
normalized coefficient triple
\[
                 (1,\alpha^2(r+1),\alpha^2r).
\]
Its symmetric functions are $\alpha,0,\alpha^2$,
so it is key $203$ in A.C2.22. Squaring the
coefficients gives key $302$. This field
automorphism transports the theta classes,
all algebraic identities below, and the
associated cyclic cover. It therefore suffices
to prove the assertion for $C_0$.

\proofparagraph{The fixed field and exact torsion generator.}
All computations take place in
\[
 F=\F_2[w]/(w^{18}+w^7+1)=\F_{4^9}.                         \tag{C2.83}
\]
An integer between $0$ and $2^{18}-1$ encodes
the polynomial whose binary digits are its
coefficients in ascending powers of $w$.
The binary modulus is $262273$, the primitive
generator is $2$, and $\alpha$ is $114587$.
The verifier checks the complete nonzero
power cycle of length $262143$. An independent
verifier also checks the polynomial
irreducibility criterion at exponents $18,9,6$.

The immutable data file
\path{code/A_F4_large_theta_witness.json}
contains a theta generator
\[
 R=(49,185454),\qquad \eta=[R-B_\Theta],
\]
whose reduced Mumford pair on (C2.80) is
\[
 U=z^2+124780z+171897,
 \qquad V=26069z+117672.                                   \tag{C2.84}
\]
The coordinates of $R$ are the original $(z,v)$
coordinates, whereas $V$ is the $Y$ coordinate
of the reduced divisor. Every displayed
coefficient is interpreted in (C2.83).
The exact checks are
\[
 [M]\eta=0,\qquad [M/p]\eta\ne0
       \quad(p=7,31,2503,4519).                             \tag{C2.85}
\]
The four factors are verified prime by trial
division, so $\eta$ has exact order $M$.

For an additional check of the primary module,
put
\[
 R_H=(1,114586),\qquad H=[4M/7]\,[R_H-B_\Theta].
\]
Its reduced pair is
\[
                   (z^2+z+1,\,114586z).
\]
The verifier proves $[7]H=0$ and
$H\ne[jM/7]\eta$ for $0\le j\le6$.
Thus it explicitly proves two independent
seven-torsion classes.
The actual $\F_4$ Frobenius polynomial from
A.C2.24 gives
\[
 \begin{aligned}
 P_4(X)&=X^4+2X^3+X^2+8X+16,\\
 \#J(F)&=68\,725\,982\,332
      =4\cdot7^2\cdot31\cdot2503\cdot4519,\\
 \#C_0(F)&=262171.
 \end{aligned}                                             \tag{C2.86}
\]
One verifier uses $\det(\pi^9-I)$ on the
integral companion matrix; two others use
integer Newton recurrences. These are the
curve and Jacobian Frobenius formulas of
\citet[II, Theorem 1.1; III, Theorem 11.1
and Corollary 11.4]{A_Milne2008}.
The group order and the independent
seven-torsion classes prove
\[
 J(F)_{\rm odd}\cong
 C_7^2\oplus C_{31}\oplus C_{2503}\oplus C_{4519}.
\]
Thus $M$ is its exact exponent. The main
certificate needs only (C2.85) and the
point-multiple identities below; this
additional argument justifies the module
structure used in the discovery search.

\proofparagraph{Divisor meaning of the finite arithmetic.}
A reduced pair $(U,V)$ means that $U$ is monic,
$\deg V<\deg U\le2$, and $U\mid V^2+hV+f$.
It represents $D(U,V)-\deg(U)W_\infty$,
where $D(U,V)$ is the finite divisor of the
ideal $(U,Y-V)$. This representation is
unique; zero is $(1,0)$
\citep[Theorem 2.2, author p.~5]{A_Lange2005}.
The checker implements the general
composition and reduction rules, including
coincident points and branch fibres.

For two pairs, compute the monic extended gcds
\[
 \begin{aligned}
 d_1&=a_0U_1+b_0U_2=\gcd(U_1,U_2),\\
 d&=e d_1+c(V_1+V_2+h)
       =\gcd(d_1,V_1+V_2+h),\qquad a=ea_0,\quad b=eb_0.
 \end{aligned}
\]
Composition is
\[
 U=\frac{U_1U_2}{d^2},\qquad
 V=\frac{aU_1V_2+bU_2V_1+c(V_1V_2+f)}d\pmod U.             \tag{C2.87}
\]
While $\deg U>2$, reduction replaces it by
\[
 U'=\frac{V^2+hV+f}U,
 \qquad V'=V+h\pmod {U'},                                  \tag{C2.88}
\]
making each new first polynomial monic.
These are the characteristic-two forms of
\citet[Algorithms 3.1--3.2, author pp.~5--6]
{A_Lange2005}, which explicitly allow the
general $h$ polynomial and even characteristic.

To recall why these are divisor identities,
the product of the two ideals is generated by
\[
 \begin{gathered}
 U_1U_2,\quad U_1(Y+V_2),\quad U_2(Y+V_1),\\
 (V_1+V_2+h)Y+(V_1V_2+f).
 \end{gathered}
\]
The Bezout relation produces $d(Y+V)$.
Locally, $d$ removes exactly the common
opposite-point pairs, giving
$I_1I_2=(d)I(U,V)$. A repeated ramification
point is removed using
$\operatorname{div}(z-z_0)=2W-2W_\infty$.
Thus (C2.87) represents the sum of classes.
For (C2.88), the norm of $Y+V$ is
$V^2+hV+f=UU'$. Its finite zero divisor is
$D(U,V)+D(U',V)$, with multiplicities, and
its only pole is at infinity. The two
degree-zero divisor classes are negatives.
Hyperelliptic conjugation changes $V$ to
$V+h$, proving the reduction rule. Monic
normalization changes only a nonzero scalar.

Every division is checked to have zero
remainder, and the defining congruence is
checked after every operation. Binary
repeated addition therefore supplies exact
identities in $\operatorname{Pic}^0(C_0)$,
not approximate numerical relations.
For a saved point $(x,y)$, the checker starts
with $([x,1],[h(x)y])$, then adds $(h,0)$.
By (C2.82) the resulting class is
$[R-B_\Theta]$.

\proofparagraph{Complete primitive orbit coverage.}
The fixed data contain $1296$ records
$(x_i,y_i,i)$ such that
\[
 \gcd(i,M)=1,
 \qquad [i]\eta=[R_i-B_\Theta]\in\Theta.                    \tag{C2.89}
\]
Every record, not a sample, is verified by
the preceding exact arithmetic. The local
orders of two at $7,31,2503,4519$ are
$3,5,1251,753$. Hence
\[
 \operatorname{ord}_M(2)=1\,570\,005,
 \qquad \varphi(M)=2\,034\,726\,480,
 \qquad \frac{\varphi(M)}{\operatorname{ord}_M(2)}=1296.     \tag{C2.90}
\]
It remains to prove that the recorded indices
belong to distinct doubling orbits.

One verification uses the primitive root
$3$ at each prime. Integer discrete logarithms
identify the unit group with
$\Z/6\oplus\Z/30\oplus\Z/2502\oplus\Z/4518$.
The lattice generated by these four period
vectors and the log vector of two has column
Hermite basis
\[
 \begin{pmatrix}
 6&0&0&2\\0&6&0&0\\0&0&6&4\\0&0&0&6
 \end{pmatrix},\qquad \det=1296.                            \tag{C2.91}
\]
Reducing from the last column to the first
gives $1296$ different residues for the
saved indices. The vector of two reduces
to zero. This proves exhaustion.

A separate verification avoids discrete-log
normalization and Hermite normal form.
Enumerate the local doubling cycles, and
label each local unit by its smallest cycle
representative and its time coordinate $t_p$.
Two global units are in one orbit exactly
when their local representatives agree and
the required time differences admit one
common shift. By the generalized Chinese
remainder theorem this is equivalent to
agreement of every pair of shifts modulo
the gcd of the two local periods. The four
local representatives and all such relative
times give $1296$ distinct signatures on
the data. A third implementation uses
different primitive roots and the equivalent
signature with local log residues modulo
$(2,6,2,6)$ and two relative times modulo
three; its $144\cdot9$ signatures also exhaust
the quotient. These are independent exact
checks of the inference from (C2.89) to
all primitive doubling orbits.

\proofparagraph{The cover and the section conclusion.}
Let $L$ represent $\eta$ over an algebraic
closure and choose $L^M\simeq\mathcal O$.
The algebra $\bigoplus_{j=0}^{M-1}L^{-j}$
defines a finite \'etale $\mu_M$-torsor
$D\to C_0$, since $M$ is odd. It is
geometrically connected because $L$ has
exact order $M$. After a finite constant
extension its deck group is cyclic of
order $M$, and its genus is $M+1$.

For every nonzero character $i$,
$H^1(C_0,L^i)$ is one dimensional. The
exact sequence
\[
 0\longrightarrow L^i\longrightarrow F_*L^{2i}
 \longrightarrow B_\Theta\otimes L^i\longrightarrow0
\]
makes its Frobenius arrow zero exactly
when $L^i\in\Theta$, as in A.C2.5.
Every primitive cycle contains a zero arrow
by (C2.89)--(C2.91). Frobenius is therefore
nilpotent on the entire primitive
conductor-$M$ summand, and that Jacobian
factor has two-rank zero. Lower-conductor
characters need not vanish.

A.C2.19, or equivalently the all-genus
criterion A.C2.20, now proves A1 for $C_0$.
Four distinct prime factors are permitted:
the primitive operator has support bound
$2^4=16$, and $M>15^2$, exactly as required
in that proof. This invokes the proved
geometric divisor argument there, not an
inference from the size of the cover.

The property is geometric. For a form over
a finitely generated field, the isomorphism,
the torsion class, its trivialization and
the deck maps are defined after a finite
extension of the whole base field. The
geometric formulation in A.C2.19--A.C2.20,
together with A.3 and the descent in A.8,
applies. An original collision would persist
over that extension. Functoriality to the
compactification handles open subcurves.
Coefficient conjugation gives key $302$,
and A.C2.22 supplies the other $46$ classes.

\proofparagraph{Reproducible certificates and exact scope.}
The fixed witness file has SHA256
\begin{center}
\path{edcf6ed3d425eb72837ae72f521f7221ee91c8a53c1b005885a4c1280d530ae6}.
\end{center}
The canonical checker
\path{code/A_F4_large_theta_canonical.py}
and its adjacent JSON verify all $1296$
relations and the order statements in
$60773$ exact divisor operations. It
imports no C implementation, search code,
accelerated polynomial engine, or tailored
coset normalization. Root and the independent
reviewer each reran it successfully.
The further independent programs
\path{code/A_F4_large_theta_root_orbits.py}
and
\path{code/A_F4_large_theta_degeneration.py},
with their JSON outputs, verify the field,
point equations, group order, reindexing
and complete orbit signatures by the
distinct methods described above.
The full mathematical reviews are
\path{audit/A_F4_large_theta_degeneration_review.md}
and
\path{audit/A_F4_all_48_root_review.md}.

The discovery code and full data are
retained under
\path{code/A_F4_large_theta_search.py/.json}.
Discovery stopped when this positive
packet was complete; no negative
enumeration is needed. A.C2.24's exact
small-primary and mixed-support exclusions
remain valid, including the entire
$73$-primary obstruction. They do not
exclude the different mixed conductor
(C2.79). The fixed-prime obstruction
A.C2.14 and the abelian-tower descent
A.C2.23 are also unchanged. This proves
the stated finite moduli class theorem;
universal A1 remains unresolved.
\end{proof}


\subsection{Every genus-two curve over \texorpdfstring{$\F_4$}{F4}}

\begin{claim}{A.C2.26}{Proved}
Every smooth proper geometrically connected
genus-two curve over $\F_4$ satisfies A1.
The conclusion holds over every finite
extension, for every finitely generated
form of any of these geometric curves,
and for every hyperbolic open subcurve.

The twelve geometric two-rank-one classes
fixed by fourth-power Frobenius are
\[
 C(A,B):\quad y^2+y=Ax^3+Bx+\frac1x,
 \qquad A\in\F_4^*,\quad B\in\F_4.                         \tag{C2.92}
\]
Explicit odd theta-torsion annihilators
$5,7,11,13,19,39$ settle all of them.
Together with A.C2.25 and the two-rank-zero
argument, this covers all $64$ geometric
genus-two moduli classes fixed by
fourth-power Frobenius. Universal A1
over arbitrary finite fields remains A.30.
\end{claim}
\begin{proof}
\proofparagraph{The exact twelve-class normalization.}
Over an algebraically closed field of
characteristic two, the genus-two
hyperelliptic quotient is separable and
admits a reduced Artin--Schreier equation.
For odd pole orders $m_j$, the genus and
two-rank formulas are
\[
 2g-2=-4+\sum_j(m_j+1),\qquad
 f(C)=\#\{\text{poles}\}-1
\]
\citep[Notation 1.1, author p.~2]{A_ElkinPries2013}.
When $g=2$ and $f(C)=1$, there are two
poles, of orders three and one. Send
them to infinity and zero. Eliminating
even terms and an additive constant gives
$ax^3+bx+c/x$ with $ac\ne0$; a unique
scaling normalizes $c=1$.
These are precisely the case $(1,3)$
normal forms of
\citet[Proposition 14, author p.~14]
{A_CardonaNartPujolas2005}.

The normalized pair $(A,B)$ is a complete
geometric invariant. An isomorphism
preserves the unique hyperelliptic pencil.
The branch points have different different
exponents, four and two, so the induced
base map fixes both and is $x\mapsto ux$.
A generator change is $y\mapsto y+v(x)$.
A nonconstant $v$ gives an even leading
pole order in $v^2+v$, whereas the difference
of the reduced equations has only odd
pole orders. Thus $v$ is constant.
The simple-pole coefficient forces $u=1$,
and then $A,B$ agree. This is also the
scaling criterion in the cited proposition.

Fourth-power Frobenius fixes the class
exactly when $A^4=A$ and $B^4=B$.
There are $3\cdot4=12$ possibilities,
all represented by actual equations
over $\F_4$. Under coefficientwise
squaring the seven representatives are
\[
 (1,0),(1,1),(1,\alpha),(\alpha,0),
 (\alpha,1),(\alpha,\alpha),(\alpha,\alpha^2),
 \qquad \alpha^2+\alpha+1=0.                               \tag{C2.93}
\]
The first two are fixed, and the others
have distinct conjugates. No descent
assertion from a field of moduli is
required for this count.

\proofparagraph{The theta origin and a norm criterion.}
Let $W_0,W_\infty$ lie above the branch
values. The valuations are
\[
 \begin{array}{c|rr}
       &x&y\\\hline
 W_0&2&-1\\
 W_\infty&-2&-3
 \end{array}
\]
and the different is $2W_0+4W_\infty$.
Consequently
\[
 \operatorname{div}(dx)=2W_0,
 \qquad B_\Theta=\mathcal O(W_0),
 \qquad \Theta=\{[R-W_0]:R\in C\}.                         \tag{C2.94}
\]
Here $B_\Theta$ is the line bundle of
locally exact differentials. In particular
zero is on theta, so a nonzero witness
must be checked. In the polynomial model
$Y=xy$, one has
\[
                 Y^2+xY=Ax^5+Bx^3+x.
\]
The pair $(x,0)$ represents the two-torsion
class $[W_0-W_\infty]$, so adding it to
$[R-W_\infty]$ gives the correct theta
class $[R-W_0]$. The fixed norm proof
below is independent of the discovery
calculation with this Jacobian model.

Let $n>1$ be odd and $d=(n-1)/2$.
Suppose a finite extension contains
$t\ne0,y_0$, polynomials $a,b$ and a
nonzero scalar $c$ such that
\[
 \begin{aligned}
 &\deg a\le d,\qquad \deg b\le d-2,\\
 &x(a^2+ab)+b^2(Ax^4+Bx^2+1)=c(x+t)^n,\\
 &a(t)+b(t)y_0=0,\qquad b(0)b(t)\ne0,\\
 &y_0^2+y_0=At^3+Bt+1/t.
 \end{aligned}                                             \tag{C2.95}
\]
Put $R=(t,y_0)$ and
\[
                            q=\frac{a+by}{x^d}.
\]
At $W_0$, the term $b(0)y$ has valuation
$-1$, strictly smaller than that of $a$;
the denominator has valuation $2d$.
Thus the pole order is exactly $n$.
At infinity, $a/x^d$ has nonnegative
valuation and $by/x^d$ has valuation at
least one. There is no other pole.
Equation (C2.95) says
\[
                  q\,\iota(q)=c\frac{(x+t)^n}{x^n}.
\]
Every zero therefore lies above $x=t$.
At $R$ its numerator is zero, while at
the conjugate point it equals $b(t)\ne0$.
This fibre is unramified, so all $n$
zeros lie at $R$. We have proved the
exact divisor identity
\[
                      \operatorname{div}(q)=n(R-W_0).      \tag{C2.96}
\]
The class $[R-W_0]$ is nonzero: otherwise
a degree-one function would make the
genus-two curve rational. Its actual odd
order $e>1$ divides $n$; no proof that
$e=n$ is needed.

\proofparagraph{Seven fixed polynomial identities.}
The following table specifies every
representative in (C2.93). Integer field
elements are binary polynomial coefficients
in ascending powers of a generator $w$.
The column $m$ is the binary modulus;
$g$ is a checked primitive multiplicative
generator. The columns $\alpha,t,y_0$
use that same encoding.
\begin{center}
\begin{tabular}{crrrrrr}
\toprule
$(A,B)$&$|F|$&$m$&$g$&$\alpha$&$(t,y_0)$&$n$\\
\midrule
$(1,0)$&16&19&2&6&$(6,2)$&13\\
$(1,1)$&4&7&2&2&$(1,2)$&5\\
$(1,\alpha)$&16&19&2&6&$(7,2)$&11\\
$(\alpha,0)$&256&283&3&189&$(80,242)$&39\\
$(\alpha,1)$&16&19&2&6&$(6,4)$&19\\
$(\alpha,\alpha)$&4&7&2&2&$(1,2)$&7\\
$(\alpha,\alpha^2)$&16&19&2&6&$(7,2)$&5\\
\bottomrule
\end{tabular}
\end{center}
The three moduli are
$w^2+w+1$, $w^4+w+1$, and
$w^8+w^4+w^3+w+1$.
The complete coefficients of $a,b,c$
for all seven identities are in
\path{code/A_F4_rank_one_witnesses.json}.
For example the degree-seven row has
\[
 a=\alpha^2+x+\alpha x^2+x^3,
 \qquad b=1+\alpha x,
 \qquad c=1.
\]
The largest row has degree $39$ and
uses only $\F_{256}$.

The fixed data have SHA256
\begin{center}
\path{82dd18cfcfac963e1a4217e03a55b39a47766ee521c6601f3704e9828b417ec3}.
\end{center}
The standalone direct verifier
\path{code/A_F4_rank_one_norm_checks.py}
and its JSON output check every coefficient
of (C2.95), the point and local conditions,
complete generator cycles, and the
twelve-class conjugation coverage.
It imports no C arithmetic, search,
jet solver or Jacobian implementation.
The independent verifier
\path{code/A_F4_rank_one_degeneration.py}
and its output use freshly implemented
binary-field and polynomial arithmetic;
all seven identities pass there too.
These are positive fixed norm certificates,
independent of any exhaustive negative scan.

\proofparagraph{Primitive vanishing and all three strata.}
Direct multiplication gives
\[
                 \langle2,-1\rangle=(\Z/n)^*
       \qquad(n=5,7,11,13,19,39).                           \tag{C2.97}
\]
At $39$ the doubling order is $12$ and
negation supplies the second coset, giving
all $24$ units. Reduction of units from
$n$ to any divisor $e$ is surjective, so
(C2.97) holds at the actual order of the
nonzero class $\eta=[R-W_0]$ from (C2.96).
Theta is symmetric because
$R+\iota(R)\sim2W_0$. Thus $\eta$ and
$-\eta$ supply a theta point on every
primitive doubling orbit at $e$.
By A.C2.5 the connected cyclic \'etale
cover of actual degree $e$ has primitive
factor of two-rank zero. A.C2.19 or
A.C2.20 proves A1. There is no extraction
of a primary component that might leave
theta, and no claim about the survival
of lower-conductor characters is needed.

Squaring coefficients gives the five
conjugate rows and all twelve classes.
Finite-extension descent, finitely
generated forms and open subcurves follow
from the geometric criterion, A.3 and
A.8, as in A.C2.25. In particular all
isomorphism and cover data are first
defined after a finite extension of
the whole base field.

Every genus-two curve has two-rank
zero, one or two. Rank one is settled
above, and rank two is A.C2.25. At rank
zero, $J[2](\overline k)=0$, and repeated
doubling excludes every nonzero geometric
two-primary torsion point. A.1 turns a
section collision into a two-primary
point difference, which is therefore
zero. The Abel map is injective, giving
the conclusion in that stratum over
every finitely generated characteristic-two
field.

For the count of geometric classes,
the invariant description of
\citet[Theorem 20, author p.~19;
Theorem 22 and Corollary 23, author
p.~21]{A_CardonaNartPujolas2005}
identifies the moduli with three invariant
coordinates. Over $\F_4$ the three
strata have respectively $4$, $12$, and
$48$ geometric classes, for a total of
$64$. This distinguishes geometric
classes from their arithmetic twists;
all twists are included in the forms
assertion. The cited classification
does not assert the prime-to-two section
theorem proved here by the certificates.

The full mathematical reviews are
\path{audit/A_F4_rank_one_degeneration_review.md}
and
\path{audit/A_F4_all_genus_two_root_review.md}.
The earlier theta-five and theta-thirteen
results remain valid and have direct
witnesses in the table. The conclusion
is the stated finite-field special case;
no universal existence of a suitable
primitive-zero cover is inferred.
\end{proof}

